\makeatletter
\newif\ifBookIndexing
\BookIndexingfalse
\newcommand{\BookBoldPage}[1]{\textbf{\hyperpage{#1}}}
\DeclareRobustCommand{\BookRecordIndex}[1]{%
  \ifBookIndexing
    \expandafter\index\expandafter{\csname BookIndex@#1\endcsname|hyperpage}%
  \fi}
\DeclareRobustCommand{\BookIndexDefinition}[1]{%
  \ifBookIndexing
    \expandafter\index\expandafter{\csname BookIndex@#1\endcsname|BookBoldPage}%
  \fi}
\DeclareRobustCommand{\BookTerm}[2]{%
  \BookRecordIndex{#1}%
  \begingroup\let\BookRecordIndex\@gobble #2\endgroup}
\AddToHook{cmd/@outputpage/before}{\let\BookRecordIndex\@gobble}
\pdfstringdefDisableCommands{%
  \def\BookTerm#1#2{#2}%
  \let\BookRecordIndex\@gobble
  \let\BookIndexDefinition\@gobble}
\expandafter\def\csname BookIndex@T0001\endcsname{liuxing@流形（manifold）}
\expandafter\let\csname BookOriginal@trManifold\endcsname\trManifold
\renewcommand{\trManifold}{\BookTerm{T0001}{\csname BookOriginal@trManifold\endcsname}}
\expandafter\def\csname BookIndex@T0002\endcsname{daibianliuxing@带边流形（manifold with boundary）}
\expandafter\let\csname BookOriginal@trManifoldWithBoundary\endcsname\trManifoldWithBoundary
\renewcommand{\trManifoldWithBoundary}{\BookTerm{T0002}{\csname BookOriginal@trManifoldWithBoundary\endcsname}}
\expandafter\def\csname BookIndex@T0003\endcsname{daijiaobian@带角边（cornered boundary）}
\expandafter\let\csname BookOriginal@trCorneredBoundary\endcsname\trCorneredBoundary
\renewcommand{\trCorneredBoundary}{\BookTerm{T0003}{\csname BookOriginal@trCorneredBoundary\endcsname}}
\expandafter\def\csname BookIndex@T0004\endcsname{jiaodian@角点（corner point）}
\expandafter\let\csname BookOriginal@trCornerPoint\endcsname\trCornerPoint
\renewcommand{\trCornerPoint}{\BookTerm{T0004}{\csname BookOriginal@trCornerPoint\endcsname}}
\expandafter\def\csname BookIndex@T0005\endcsname{yuweishu@余维数（codimension）}
\expandafter\let\csname BookOriginal@trCodimension\endcsname\trCodimension
\renewcommand{\trCodimension}{\BookTerm{T0005}{\csname BookOriginal@trCodimension\endcsname}}
\expandafter\def\csname BookIndex@T0006\endcsname{fenceng@分层（stratum）}
\expandafter\let\csname BookOriginal@trStratum\endcsname\trStratum
\renewcommand{\trStratum}{\BookTerm{T0006}{\csname BookOriginal@trStratum\endcsname}}
\expandafter\def\csname BookIndex@T0007\endcsname{bianjiefenceng@边界分层（boundary stratum）}
\expandafter\let\csname BookOriginal@trBoundaryStratum\endcsname\trBoundaryStratum
\renewcommand{\trBoundaryStratum}{\BookTerm{T0007}{\csname BookOriginal@trBoundaryStratum\endcsname}}
\expandafter\def\csname BookIndex@T0008\endcsname{bianjiemian@边界面（boundary face）}
\expandafter\let\csname BookOriginal@trBoundaryFace\endcsname\trBoundaryFace
\renewcommand{\trBoundaryFace}{\BookTerm{T0008}{\csname BookOriginal@trBoundaryFace\endcsname}}
\expandafter\def\csname BookIndex@T0009\endcsname{zixiangjiao@自相交（self intersection）}
\expandafter\let\csname BookOriginal@trSelfIntersection\endcsname\trSelfIntersection
\renewcommand{\trSelfIntersection}{\BookTerm{T0009}{\csname BookOriginal@trSelfIntersection\endcsname}}
\expandafter\def\csname BookIndex@T0010\endcsname{jingxiang@径向（radial）}
\expandafter\let\csname BookOriginal@trRadial\endcsname\trRadial
\renewcommand{\trRadial}{\BookTerm{T0010}{\csname BookOriginal@trRadial\endcsname}}
\expandafter\def\csname BookIndex@T0011\endcsname{ya@芽（germ）}
\expandafter\let\csname BookOriginal@trGerm\endcsname\trGerm
\renewcommand{\trGerm}{\BookTerm{T0011}{\csname BookOriginal@trGerm\endcsname}}
\expandafter\def\csname BookIndex@T0012\endcsname{jinzhiji@紧支集（compactly supported）}
\expandafter\let\csname BookOriginal@trCompactlySupported\endcsname\trCompactlySupported
\renewcommand{\trCompactlySupported}{\BookTerm{T0012}{\csname BookOriginal@trCompactlySupported\endcsname}}
\expandafter\def\csname BookIndex@T0013\endcsname{tapuqun@拓扑群（topological group）}
\expandafter\let\csname BookOriginal@trTopologicalGroup\endcsname\trTopologicalGroup
\renewcommand{\trTopologicalGroup}{\BookTerm{T0013}{\csname BookOriginal@trTopologicalGroup\endcsname}}
\expandafter\def\csname BookIndex@T0014\endcsname{jubuchangzhi@局部常值（locally constant）}
\expandafter\let\csname BookOriginal@trLocallyConstant\endcsname\trLocallyConstant
\renewcommand{\trLocallyConstant}{\BookTerm{T0014}{\csname BookOriginal@trLocallyConstant\endcsname}}
\expandafter\def\csname BookIndex@T0015\endcsname{ziliuxing@子流形（submanifold）}
\expandafter\let\csname BookOriginal@trSubmanifold\endcsname\trSubmanifold
\renewcommand{\trSubmanifold}{\BookTerm{T0015}{\csname BookOriginal@trSubmanifold\endcsname}}
\expandafter\def\csname BookIndex@T0016\endcsname{chaoqumian@超曲面（hypersurface）}
\expandafter\let\csname BookOriginal@trHypersurface\endcsname\trHypersurface
\renewcommand{\trHypersurface}{\BookTerm{T0016}{\csname BookOriginal@trHypersurface\endcsname}}
\expandafter\def\csname BookIndex@T0017\endcsname{tapu@拓扑（topology）}
\expandafter\let\csname BookOriginal@trTopology\endcsname\trTopology
\renewcommand{\trTopology}{\BookTerm{T0017}{\csname BookOriginal@trTopology\endcsname}}
\expandafter\def\csname BookIndex@T0018\endcsname{qiumian@球面（sphere）}
\expandafter\let\csname BookOriginal@trSphere\endcsname\trSphere
\renewcommand{\trSphere}{\BookTerm{T0018}{\csname BookOriginal@trSphere\endcsname}}
\expandafter\def\csname BookIndex@T0019\endcsname{yuanzhou@圆周（circle）}
\expandafter\let\csname BookOriginal@trCircle\endcsname\trCircle
\renewcommand{\trCircle}{\BookTerm{T0019}{\csname BookOriginal@trCircle\endcsname}}
\expandafter\def\csname BookIndex@T0020\endcsname{huanmian@环面（torus）}
\expandafter\let\csname BookOriginal@trTorus\endcsname\trTorus
\renewcommand{\trTorus}{\BookTerm{T0020}{\csname BookOriginal@trTorus\endcsname}}
\expandafter\def\csname BookIndex@T0021\endcsname{yuanpan@圆盘（disk）}
\expandafter\let\csname BookOriginal@trDisk\endcsname\trDisk
\renewcommand{\trDisk}{\BookTerm{T0021}{\csname BookOriginal@trDisk\endcsname}}
\expandafter\def\csname BookIndex@T0022\endcsname{qumian@曲面（surface）}
\expandafter\let\csname BookOriginal@trSurface\endcsname\trSurface
\renewcommand{\trSurface}{\BookTerm{T0022}{\csname BookOriginal@trSurface\endcsname}}
\expandafter\def\csname BookIndex@T0023\endcsname{guanghua@光滑（smooth）}
\expandafter\let\csname BookOriginal@trSmooth\endcsname\trSmooth
\renewcommand{\trSmooth}{\BookTerm{T0023}{\csname BookOriginal@trSmooth\endcsname}}
\expandafter\def\csname BookIndex@T0024\endcsname{zuobiaoka@坐标卡（chart）}
\expandafter\let\csname BookOriginal@trChart\endcsname\trChart
\renewcommand{\trChart}{\BookTerm{T0024}{\csname BookOriginal@trChart\endcsname}}
\expandafter\def\csname BookIndex@T0025\endcsname{linjiedian@临界点（critical point）}
\expandafter\let\csname BookOriginal@trCriticalPoint\endcsname\trCriticalPoint
\renewcommand{\trCriticalPoint}{\BookTerm{T0025}{\csname BookOriginal@trCriticalPoint\endcsname}}
\expandafter\def\csname BookIndex@T0026\endcsname{linjiezhi@临界值（critical value）}
\expandafter\let\csname BookOriginal@trCriticalValue\endcsname\trCriticalValue
\renewcommand{\trCriticalValue}{\BookTerm{T0026}{\csname BookOriginal@trCriticalValue\endcsname}}
\expandafter\def\csname BookIndex@T0027\endcsname{zhengzezhi@正则值（regular value）}
\expandafter\let\csname BookOriginal@trRegularValue\endcsname\trRegularValue
\renewcommand{\trRegularValue}{\BookTerm{T0027}{\csname BookOriginal@trRegularValue\endcsname}}
\expandafter\def\csname BookIndex@T0028\endcsname{zhengze@正则（regular）}
\expandafter\let\csname BookOriginal@trRegular\endcsname\trRegular
\renewcommand{\trRegular}{\BookTerm{T0028}{\csname BookOriginal@trRegular\endcsname}}
\expandafter\def\csname BookIndex@T0029\endcsname{shuipingji@水平集（level set）}
\expandafter\let\csname BookOriginal@trLevelSet\endcsname\trLevelSet
\renewcommand{\trLevelSet}{\BookTerm{T0029}{\csname BookOriginal@trLevelSet\endcsname}}
\expandafter\def\csname BookIndex@T0030\endcsname{cishuipingji@次水平集（sublevel set）}
\expandafter\let\csname BookOriginal@trSublevelSet\endcsname\trSublevelSet
\renewcommand{\trSublevelSet}{\BookTerm{T0030}{\csname BookOriginal@trSublevelSet\endcsname}}
\expandafter\def\csname BookIndex@T0031\endcsname{duijing@对径（antipodal）}
\expandafter\let\csname BookOriginal@trAntipodal\endcsname\trAntipodal
\renewcommand{\trAntipodal}{\BookTerm{T0031}{\csname BookOriginal@trAntipodal\endcsname}}
\expandafter\def\csname BookIndex@T0032\endcsname{qipao@气泡（bubble）}
\expandafter\let\csname BookOriginal@trBubble\endcsname\trBubble
\renewcommand{\trBubble}{\BookTerm{T0032}{\csname BookOriginal@trBubble\endcsname}}
\expandafter\def\csname BookIndex@T0033\endcsname{qidian@奇点（singular point）}
\expandafter\let\csname BookOriginal@trSingularPoint\endcsname\trSingularPoint
\renewcommand{\trSingularPoint}{\BookTerm{T0033}{\csname BookOriginal@trSingularPoint\endcsname}}
\expandafter\def\csname BookIndex@T0034\endcsname{tuihua@退化（degenerate）}
\expandafter\let\csname BookOriginal@trDegenerate\endcsname\trDegenerate
\renewcommand{\trDegenerate}{\BookTerm{T0034}{\csname BookOriginal@trDegenerate\endcsname}}
\expandafter\def\csname BookIndex@T0035\endcsname{feituihua@非退化（nondegenerate）}
\expandafter\let\csname BookOriginal@trNondegenerate\endcsname\trNondegenerate
\renewcommand{\trNondegenerate}{\BookTerm{T0035}{\csname BookOriginal@trNondegenerate\endcsname}}
\expandafter\def\csname BookIndex@T0036\endcsname{fuding@负定（negative definite）}
\expandafter\let\csname BookOriginal@trNegativeDefinite\endcsname\trNegativeDefinite
\renewcommand{\trNegativeDefinite}{\BookTerm{T0036}{\csname BookOriginal@trNegativeDefinite\endcsname}}
\expandafter\def\csname BookIndex@T0037\endcsname{guli@孤立（isolated）}
\expandafter\let\csname BookOriginal@trIsolated\endcsname\trIsolated
\renewcommand{\trIsolated}{\BookTerm{T0037}{\csname BookOriginal@trIsolated\endcsname}}
\expandafter\def\csname BookIndex@T0038\endcsname{sheying@射影（projective）}
\expandafter\let\csname BookOriginal@trProjective\endcsname\trProjective
\renewcommand{\trProjective}{\BookTerm{T0038}{\csname BookOriginal@trProjective\endcsname}}
\expandafter\def\csname BookIndex@T0039\endcsname{sheyingkongjian@射影空间（projective space）}
\expandafter\let\csname BookOriginal@trProjectiveSpace\endcsname\trProjectiveSpace
\renewcommand{\trProjectiveSpace}{\BookTerm{T0039}{\csname BookOriginal@trProjectiveSpace\endcsname}}
\expandafter\def\csname BookIndex@T0040\endcsname{qicizuobiao@齐次坐标（homogeneous coordinates）}
\expandafter\let\csname BookOriginal@trHomogeneousCoordinates\endcsname\trHomogeneousCoordinates
\renewcommand{\trHomogeneousCoordinates}{\BookTerm{T0040}{\csname BookOriginal@trHomogeneousCoordinates\endcsname}}
\expandafter\def\csname BookIndex@T0041\endcsname{haisai@海塞（hessian）}
\expandafter\let\csname BookOriginal@trHessian\endcsname\trHessian
\renewcommand{\trHessian}{\BookTerm{T0041}{\csname BookOriginal@trHessian\endcsname}}
\expandafter\def\csname BookIndex@T0042\endcsname{ercixing@二次型（quadratic form）}
\expandafter\let\csname BookOriginal@trQuadraticForm\endcsname\trQuadraticForm
\renewcommand{\trQuadraticForm}{\BookTerm{T0042}{\csname BookOriginal@trQuadraticForm\endcsname}}
\expandafter\def\csname BookIndex@T0043\endcsname{liangdingyi@良定义（well defined）}
\expandafter\let\csname BookOriginal@trWellDefined\endcsname\trWellDefined
\renewcommand{\trWellDefined}{\BookTerm{T0043}{\csname BookOriginal@trWellDefined\endcsname}}
\expandafter\def\csname BookIndex@T0044\endcsname{moersihanshu@莫尔斯函数（Morse function）}
\expandafter\let\csname BookOriginal@trMorseFunction\endcsname\trMorseFunction
\renewcommand{\trMorseFunction}{\BookTerm{T0044}{\csname BookOriginal@trMorseFunction\endcsname}}
\expandafter\def\csname BookIndex@T0045\endcsname{moersiyinli@莫尔斯引理（Morse lemma）}
\expandafter\let\csname BookOriginal@trMorseLemma\endcsname\trMorseLemma
\renewcommand{\trMorseLemma}{\BookTerm{T0045}{\csname BookOriginal@trMorseLemma\endcsname}}
\expandafter\def\csname BookIndex@T0046\endcsname{moersililun@莫尔斯理论（Morse theory）}
\expandafter\let\csname BookOriginal@trMorseTheory\endcsname\trMorseTheory
\renewcommand{\trMorseTheory}{\BookTerm{T0046}{\csname BookOriginal@trMorseTheory\endcsname}}
\expandafter\def\csname BookIndex@T0047\endcsname{tidu@梯度（gradient）}
\expandafter\let\csname BookOriginal@trGradient\endcsname\trGradient
\renewcommand{\trGradient}{\BookTerm{T0047}{\csname BookOriginal@trGradient\endcsname}}
\expandafter\def\csname BookIndex@T0048\endcsname{xiangliangchang@向量场（vector field）}
\expandafter\let\csname BookOriginal@trVectorField\endcsname\trVectorField
\renewcommand{\trVectorField}{\BookTerm{T0048}{\csname BookOriginal@trVectorField\endcsname}}
\expandafter\def\csname BookIndex@T0049\endcsname{xiangliangkongjian@向量空间（vector space）}
\expandafter\let\csname BookOriginal@trVectorSpace\endcsname\trVectorSpace
\renewcommand{\trVectorSpace}{\BookTerm{T0049}{\csname BookOriginal@trVectorSpace\endcsname}}
\expandafter\def\csname BookIndex@T0050\endcsname{shipei@适配（adapted）}
\expandafter\let\csname BookOriginal@trAdapted\endcsname\trAdapted
\renewcommand{\trAdapted}{\BookTerm{T0050}{\csname BookOriginal@trAdapted\endcsname}}
\expandafter\def\csname BookIndex@T0051\endcsname{neixiang@内向（incoming）}
\expandafter\let\csname BookOriginal@trIncoming\endcsname\trIncoming
\renewcommand{\trIncoming}{\BookTerm{T0051}{\csname BookOriginal@trIncoming\endcsname}}
\expandafter\def\csname BookIndex@T0052\endcsname{lingquan@领圈（collar）}
\expandafter\let\csname BookOriginal@trCollar\endcsname\trCollar
\renewcommand{\trCollar}{\BookTerm{T0052}{\csname BookOriginal@trCollar\endcsname}}
\expandafter\def\csname BookIndex@T0053\endcsname{zhuding@驻定（stationary）}
\expandafter\let\csname BookOriginal@trStationary\endcsname\trStationary
\renewcommand{\trStationary}{\BookTerm{T0053}{\csname BookOriginal@trStationary\endcsname}}
\expandafter\def\csname BookIndex@T0054\endcsname{guidao@轨道（trajectory）}
\expandafter\let\csname BookOriginal@trTrajectory\endcsname\trTrajectory
\renewcommand{\trTrajectory}{\BookTerm{T0054}{\csname BookOriginal@trTrajectory\endcsname}}
\expandafter\def\csname BookIndex@T0055\endcsname{jifenquxian@积分曲线（integral curve）}
\expandafter\let\csname BookOriginal@trIntegralCurve\endcsname\trIntegralCurve
\renewcommand{\trIntegralCurve}{\BookTerm{T0055}{\csname BookOriginal@trIntegralCurve\endcsname}}
\expandafter\def\csname BookIndex@T0056\endcsname{liu@流（flow）}
\expandafter\let\csname BookOriginal@trFlow\endcsname\trFlow
\renewcommand{\trFlow}{\BookTerm{T0056}{\csname BookOriginal@trFlow\endcsname}}
\expandafter\def\csname BookIndex@T0057\endcsname{wendingliuxing@稳定流形（stable manifold）}
\expandafter\let\csname BookOriginal@trStableManifold\endcsname\trStableManifold
\renewcommand{\trStableManifold}{\BookTerm{T0057}{\csname BookOriginal@trStableManifold\endcsname}}
\expandafter\def\csname BookIndex@T0058\endcsname{buwendingliuxing@不稳定流形（unstable manifold）}
\expandafter\let\csname BookOriginal@trUnstableManifold\endcsname\trUnstableManifold
\renewcommand{\trUnstableManifold}{\BookTerm{T0058}{\csname BookOriginal@trUnstableManifold\endcsname}}
\expandafter\def\csname BookIndex@T0059\endcsname{zhibiao@指标（index）}
\expandafter\let\csname BookOriginal@trIndex\endcsname\trIndex
\renewcommand{\trIndex}{\BookTerm{T0059}{\csname BookOriginal@trIndex\endcsname}}
\expandafter\def\csname BookIndex@T0060\endcsname{jinzhi@紧致（compact）}
\expandafter\let\csname BookOriginal@trCompact\endcsname\trCompact
\renewcommand{\trCompact}{\BookTerm{T0060}{\csname BookOriginal@trCompact\endcsname}}
\expandafter\def\csname BookIndex@T0061\endcsname{jinhua@紧化（compactification）}
\expandafter\let\csname BookOriginal@trCompactification\endcsname\trCompactification
\renewcommand{\trCompactification}{\BookTerm{T0061}{\csname BookOriginal@trCompactification\endcsname}}
\expandafter\def\csname BookIndex@T0062\endcsname{liantong@连通（connected）}
\expandafter\let\csname BookOriginal@trConnected\endcsname\trConnected
\renewcommand{\trConnected}{\BookTerm{T0062}{\csname BookOriginal@trConnected\endcsname}}
\expandafter\def\csname BookIndex@T0063\endcsname{liantongfenzhi@连通分支（connected component）}
\expandafter\let\csname BookOriginal@trConnectedComponent\endcsname\trConnectedComponent
\renewcommand{\trConnectedComponent}{\BookTerm{T0063}{\csname BookOriginal@trConnectedComponent\endcsname}}
\expandafter\def\csname BookIndex@T0064\endcsname{danliantong@单连通（simply connected）}
\expandafter\let\csname BookOriginal@trSimplyConnected\endcsname\trSimplyConnected
\renewcommand{\trSimplyConnected}{\BookTerm{T0064}{\csname BookOriginal@trSimplyConnected\endcsname}}
\expandafter\def\csname BookIndex@T0065\endcsname{choumi@稠密（dense）}
\expandafter\let\csname BookOriginal@trDense\endcsname\trDense
\renewcommand{\trDense}{\BookTerm{T0065}{\csname BookOriginal@trDense\endcsname}}
\expandafter\def\csname BookIndex@T0066\endcsname{kai@开（open）}
\expandafter\let\csname BookOriginal@trOpen\endcsname\trOpen
\renewcommand{\trOpen}{\BookTerm{T0066}{\csname BookOriginal@trOpen\endcsname}}
\expandafter\def\csname BookIndex@T0067\endcsname{tongyou@通有（generic）}
\expandafter\let\csname BookOriginal@trGeneric\endcsname\trGeneric
\renewcommand{\trGeneric}{\BookTerm{T0067}{\csname BookOriginal@trGeneric\endcsname}}
\expandafter\def\csname BookIndex@T0068\endcsname{tongyouxing@通有性（genericity）}
\expandafter\let\csname BookOriginal@trGenericity\endcsname\trGenericity
\renewcommand{\trGenericity}{\BookTerm{T0068}{\csname BookOriginal@trGenericity\endcsname}}
\expandafter\def\csname BookIndex@T0069\endcsname{hengjiexing@横截性（transversality）}
\expandafter\let\csname BookOriginal@trTransversality\endcsname\trTransversality
\renewcommand{\trTransversality}{\BookTerm{T0069}{\csname BookOriginal@trTransversality\endcsname}}
\expandafter\def\csname BookIndex@T0070\endcsname{hengjie@横截（transverse）}
\expandafter\let\csname BookOriginal@trTransverse\endcsname\trTransverse
\renewcommand{\trTransverse}{\BookTerm{T0070}{\csname BookOriginal@trTransverse\endcsname}}
\expandafter\def\csname BookIndex@T0071\endcsname{jinmei@浸没（submersion）}
\expandafter\let\csname BookOriginal@trSubmersion\endcsname\trSubmersion
\renewcommand{\trSubmersion}{\BookTerm{T0071}{\csname BookOriginal@trSubmersion\endcsname}}
\expandafter\def\csname BookIndex@T0072\endcsname{jinru@浸入（immersion）}
\expandafter\let\csname BookOriginal@trImmersion\endcsname\trImmersion
\renewcommand{\trImmersion}{\BookTerm{T0072}{\csname BookOriginal@trImmersion\endcsname}}
\expandafter\def\csname BookIndex@T0073\endcsname{baohanyingshe@包含映射（inclusion map）}
\expandafter\let\csname BookOriginal@trInclusionMap\endcsname\trInclusionMap
\renewcommand{\trInclusionMap}{\BookTerm{T0073}{\csname BookOriginal@trInclusionMap\endcsname}}
\expandafter\def\csname BookIndex@T0074\endcsname{touying@投影（projection）}
\expandafter\let\csname BookOriginal@trProjection\endcsname\trProjection
\renewcommand{\trProjection}{\BookTerm{T0074}{\csname BookOriginal@trProjection\endcsname}}
\expandafter\def\csname BookIndex@T0075\endcsname{qianru@嵌入（embedding）}
\expandafter\let\csname BookOriginal@trEmbedding\endcsname\trEmbedding
\renewcommand{\trEmbedding}{\BookTerm{T0075}{\csname BookOriginal@trEmbedding\endcsname}}
\expandafter\def\csname BookIndex@T0076\endcsname{canshuhua@参数化（parametrization）}
\expandafter\let\csname BookOriginal@trParametrization\endcsname\trParametrization
\renewcommand{\trParametrization}{\BookTerm{T0076}{\csname BookOriginal@trParametrization\endcsname}}
\expandafter\def\csname BookIndex@T0077\endcsname{weifentongpei@微分同胚（diffeomorphism）}
\expandafter\let\csname BookOriginal@trDiffeomorphism\endcsname\trDiffeomorphism
\renewcommand{\trDiffeomorphism}{\BookTerm{T0077}{\csname BookOriginal@trDiffeomorphism\endcsname}}
\expandafter\def\csname BookIndex@T0078\endcsname{tongpei@同胚（homeomorphism）}
\expandafter\let\csname BookOriginal@trHomeomorphism\endcsname\trHomeomorphism
\renewcommand{\trHomeomorphism}{\BookTerm{T0078}{\csname BookOriginal@trHomeomorphism\endcsname}}
\expandafter\def\csname BookIndex@T0079\endcsname{qiekongjian@切空间（tangent space）}
\expandafter\let\csname BookOriginal@trTangentSpace\endcsname\trTangentSpace
\renewcommand{\trTangentSpace}{\BookTerm{T0079}{\csname BookOriginal@trTangentSpace\endcsname}}
\expandafter\def\csname BookIndex@T0080\endcsname{qiecong@切丛（tangent bundle）}
\expandafter\let\csname BookOriginal@trTangentBundle\endcsname\trTangentBundle
\renewcommand{\trTangentBundle}{\BookTerm{T0080}{\csname BookOriginal@trTangentBundle\endcsname}}
\expandafter\def\csname BookIndex@T0081\endcsname{yuqie@余切（cotangent）}
\expandafter\let\csname BookOriginal@trCotangent\endcsname\trCotangent
\renewcommand{\trCotangent}{\BookTerm{T0081}{\csname BookOriginal@trCotangent\endcsname}}
\expandafter\def\csname BookIndex@T0082\endcsname{xiangliangcong@向量丛（vector bundle）}
\expandafter\let\csname BookOriginal@trVectorBundle\endcsname\trVectorBundle
\renewcommand{\trVectorBundle}{\BookTerm{T0082}{\csname BookOriginal@trVectorBundle\endcsname}}
\expandafter\def\csname BookIndex@T0083\endcsname{xianweicong@纤维丛（fiber bundle）}
\expandafter\let\csname BookOriginal@trFiberBundle\endcsname\trFiberBundle
\renewcommand{\trFiberBundle}{\BookTerm{T0083}{\csname BookOriginal@trFiberBundle\endcsname}}
\expandafter\def\csname BookIndex@T0084\endcsname{facong@法丛（normal bundle）}
\expandafter\let\csname BookOriginal@trNormalBundle\endcsname\trNormalBundle
\renewcommand{\trNormalBundle}{\BookTerm{T0084}{\csname BookOriginal@trNormalBundle\endcsname}}
\expandafter\def\csname BookIndex@T0085\endcsname{faxiangliang@法向量（normal vector）}
\expandafter\let\csname BookOriginal@trNormalVector\endcsname\trNormalVector
\renewcommand{\trNormalVector}{\BookTerm{T0085}{\csname BookOriginal@trNormalVector\endcsname}}
\expandafter\def\csname BookIndex@T0086\endcsname{guanzhuanglinyu@管状邻域（tubular neighborhood）}
\expandafter\let\csname BookOriginal@trTubularNeighborhood\endcsname\trTubularNeighborhood
\renewcommand{\trTubularNeighborhood}{\BookTerm{T0086}{\csname BookOriginal@trTubularNeighborhood\endcsname}}
\expandafter\def\csname BookIndex@T0087\endcsname{xianwei@纤维（fiber）}
\expandafter\let\csname BookOriginal@trFiber\endcsname\trFiber
\renewcommand{\trFiber}{\BookTerm{T0087}{\csname BookOriginal@trFiber\endcsname}}
\expandafter\def\csname BookIndex@T0088\endcsname{jiemian@截面（section）}
\expandafter\let\csname BookOriginal@trSection\endcsname\trSection
\renewcommand{\trSection}{\BookTerm{T0088}{\csname BookOriginal@trSection\endcsname}}
\expandafter\def\csname BookIndex@T0089\endcsname{lingjiemian@零截面（zero section）}
\expandafter\let\csname BookOriginal@trZeroSection\endcsname\trZeroSection
\renewcommand{\trZeroSection}{\BookTerm{T0089}{\csname BookOriginal@trZeroSection\endcsname}}
\expandafter\def\csname BookIndex@T0090\endcsname{pingfanhua@平凡化（trivialization）}
\expandafter\let\csname BookOriginal@trTrivialization\endcsname\trTrivialization
\renewcommand{\trTrivialization}{\BookTerm{T0090}{\csname BookOriginal@trTrivialization\endcsname}}
\expandafter\def\csname BookIndex@T0091\endcsname{duliang@度量（metric）}
\expandafter\let\csname BookOriginal@trMetric\endcsname\trMetric
\renewcommand{\trMetric}{\BookTerm{T0091}{\csname BookOriginal@trMetric\endcsname}}
\expandafter\def\csname BookIndex@T0092\endcsname{zhishushuaijian@指数衰减（exponential decay）}
\expandafter\let\csname BookOriginal@trExponentialDecay\endcsname\trExponentialDecay
\renewcommand{\trExponentialDecay}{\BookTerm{T0092}{\csname BookOriginal@trExponentialDecay\endcsname}}
\expandafter\def\csname BookIndex@T0093\endcsname{danweifenjie@单位分解（partition of unity）}
\expandafter\let\csname BookOriginal@trPartitionOfUnity\endcsname\trPartitionOfUnity
\renewcommand{\trPartitionOfUnity}{\BookTerm{T0093}{\csname BookOriginal@trPartitionOfUnity\endcsname}}
\expandafter\def\csname BookIndex@T0094\endcsname{zhiji@支集（support）}
\expandafter\let\csname BookOriginal@trSupport\endcsname\trSupport
\renewcommand{\trSupport}{\BookTerm{T0094}{\csname BookOriginal@trSupport\endcsname}}
\expandafter\def\csname BookIndex@T0095\endcsname{shoubing@手柄（handle）}
\expandafter\let\csname BookOriginal@trHandle\endcsname\trHandle
\renewcommand{\trHandle}{\BookTerm{T0095}{\csname BookOriginal@trHandle\endcsname}}
\expandafter\def\csname BookIndex@T0096\endcsname{fuxing@复形（complex）}
\expandafter\let\csname BookOriginal@trComplex\endcsname\trComplex
\renewcommand{\trComplex}{\BookTerm{T0096}{\csname BookOriginal@trComplex\endcsname}}
\expandafter\def\csname BookIndex@T0097\endcsname{shengchengyuan@生成元（generator）}
\expandafter\let\csname BookOriginal@trGenerator\endcsname\trGenerator
\renewcommand{\trGenerator}{\BookTerm{T0097}{\csname BookOriginal@trGenerator\endcsname}}
\expandafter\def\csname BookIndex@T0098\endcsname{bujiaobing@不交并（disjoint union）}
\expandafter\let\csname BookOriginal@trDisjointUnion\endcsname\trDisjointUnion
\renewcommand{\trDisjointUnion}{\BookTerm{T0098}{\csname BookOriginal@trDisjointUnion\endcsname}}
\expandafter\def\csname BookIndex@T0099\endcsname{abeierqun@阿贝尔群（abelian group）}
\expandafter\let\csname BookOriginal@trAbelianGroup\endcsname\trAbelianGroup
\renewcommand{\trAbelianGroup}{\BookTerm{T0099}{\csname BookOriginal@trAbelianGroup\endcsname}}
\expandafter\def\csname BookIndex@T0100\endcsname{cishu@次数（degree）}
\expandafter\let\csname BookOriginal@trDegree\endcsname\trDegree
\renewcommand{\trDegree}{\BookTerm{T0100}{\csname BookOriginal@trDegree\endcsname}}
\expandafter\def\csname BookIndex@T0101\endcsname{baoqiang@胞腔（cellular）}
\expandafter\let\csname BookOriginal@trCellular\endcsname\trCellular
\renewcommand{\trCellular}{\BookTerm{T0101}{\csname BookOriginal@trCellular\endcsname}}
\expandafter\def\csname BookIndex@T0102\endcsname{baoqiangfenjie@胞腔分解（cellular decomposition）}
\expandafter\let\csname BookOriginal@trCellularDecomposition\endcsname\trCellularDecomposition
\renewcommand{\trCellularDecomposition}{\BookTerm{T0102}{\csname BookOriginal@trCellularDecomposition\endcsname}}
\expandafter\def\csname BookIndex@T0103\endcsname{cedexian@测地线（geodesic）}
\expandafter\let\csname BookOriginal@trGeodesic\endcsname\trGeodesic
\renewcommand{\trGeodesic}{\BookTerm{T0103}{\csname BookOriginal@trGeodesic\endcsname}}
\expandafter\def\csname BookIndex@T0104\endcsname{moersitongdiao@莫尔斯同调（Morse homology）}
\expandafter\let\csname BookOriginal@trMorseHomology\endcsname\trMorseHomology
\renewcommand{\trMorseHomology}{\BookTerm{T0104}{\csname BookOriginal@trMorseHomology\endcsname}}
\expandafter\def\csname BookIndex@T0105\endcsname{mokongjian@模空间（moduli space）}
\expandafter\let\csname BookOriginal@trModuliSpace\endcsname\trModuliSpace
\renewcommand{\trModuliSpace}{\BookTerm{T0105}{\csname BookOriginal@trModuliSpace\endcsname}}
\expandafter\def\csname BookIndex@T0106\endcsname{xujibenlian@虚基本链（virtual fundamental chain）}
\expandafter\let\csname BookOriginal@trVirtualFundamentalChain\endcsname\trVirtualFundamentalChain
\renewcommand{\trVirtualFundamentalChain}{\BookTerm{T0106}{\csname BookOriginal@trVirtualFundamentalChain\endcsname}}
\expandafter\def\csname BookIndex@T0107\endcsname{zhanhe@粘合（gluing）}
\expandafter\let\csname BookOriginal@trGluing\endcsname\trGluing
\renewcommand{\trGluing}{\BookTerm{T0107}{\csname BookOriginal@trGluing\endcsname}}
\expandafter\def\csname BookIndex@T0108\endcsname{pingxingyidong@平行移动（parallel transport）}
\expandafter\let\csname BookOriginal@trParallelTransport\endcsname\trParallelTransport
\renewcommand{\trParallelTransport}{\BookTerm{T0108}{\csname BookOriginal@trParallelTransport\endcsname}}
\expandafter\def\csname BookIndex@T0109\endcsname{zicong@子丛（subbundle）}
\expandafter\let\csname BookOriginal@trSubbundle\endcsname\trSubbundle
\renewcommand{\trSubbundle}{\BookTerm{T0109}{\csname BookOriginal@trSubbundle\endcsname}}
\expandafter\def\csname BookIndex@T0110\endcsname{nijin@逆紧（proper）}
\expandafter\let\csname BookOriginal@trProper\endcsname\trProper
\renewcommand{\trProper}{\BookTerm{T0110}{\csname BookOriginal@trProper\endcsname}}
\expandafter\def\csname BookIndex@T0111\endcsname{tisheng@提升（lift）}
\expandafter\let\csname BookOriginal@trLift\endcsname\trLift
\renewcommand{\trLift}{\BookTerm{T0111}{\csname BookOriginal@trLift\endcsname}}
\expandafter\def\csname BookIndex@T0112\endcsname{tongdiao@同调（homology）}
\expandafter\let\csname BookOriginal@trHomology\endcsname\trHomology
\renewcommand{\trHomology}{\BookTerm{T0112}{\csname BookOriginal@trHomology\endcsname}}
\expandafter\def\csname BookIndex@T0113\endcsname{shangtongdiao@上同调（cohomology）}
\expandafter\let\csname BookOriginal@trCohomology\endcsname\trCohomology
\renewcommand{\trCohomology}{\BookTerm{T0113}{\csname BookOriginal@trCohomology\endcsname}}
\expandafter\def\csname BookIndex@T0114\endcsname{bubianliang@不变量（invariant）}
\expandafter\let\csname BookOriginal@trInvariant\endcsname\trInvariant
\renewcommand{\trInvariant}{\BookTerm{T0114}{\csname BookOriginal@trInvariant\endcsname}}
\expandafter\def\csname BookIndex@T0115\endcsname{bubianxing@不变性（invariance）}
\expandafter\let\csname BookOriginal@trInvariance\endcsname\trInvariance
\renewcommand{\trInvariance}{\BookTerm{T0115}{\csname BookOriginal@trInvariance\endcsname}}
\expandafter\def\csname BookIndex@T0116\endcsname{qiyi@奇异（singular）}
\expandafter\let\csname BookOriginal@trSingular\endcsname\trSingular
\renewcommand{\trSingular}{\BookTerm{T0116}{\csname BookOriginal@trSingular\endcsname}}
\expandafter\def\csname BookIndex@T0117\endcsname{feiqiyi@非奇异（nonsingular）}
\expandafter\let\csname BookOriginal@trNonsingular\endcsname\trNonsingular
\renewcommand{\trNonsingular}{\BookTerm{T0117}{\csname BookOriginal@trNonsingular\endcsname}}
\expandafter\def\csname BookIndex@T0118\endcsname{tonghen@同痕（isotopy）}
\expandafter\let\csname BookOriginal@trIsotopy\endcsname\trIsotopy
\renewcommand{\trIsotopy}{\BookTerm{T0118}{\csname BookOriginal@trIsotopy\endcsname}}
\expandafter\def\csname BookIndex@T0119\endcsname{jiechu@接触（contact）}
\expandafter\let\csname BookOriginal@trContact\endcsname\trContact
\renewcommand{\trContact}{\BookTerm{T0119}{\csname BookOriginal@trContact\endcsname}}
\expandafter\def\csname BookIndex@T0120\endcsname{biaojia@标架（frame）}
\expandafter\let\csname BookOriginal@trFrame\endcsname\trFrame
\renewcommand{\trFrame}{\BookTerm{T0120}{\csname BookOriginal@trFrame\endcsname}}
\expandafter\def\csname BookIndex@T0121\endcsname{xiangdui@相对（relative）}
\expandafter\let\csname BookOriginal@trRelative\endcsname\trRelative
\renewcommand{\trRelative}{\BookTerm{T0121}{\csname BookOriginal@trRelative\endcsname}}
\expandafter\def\csname BookIndex@T0122\endcsname{lianfuxing@链复形（chain complex）}
\expandafter\let\csname BookOriginal@trChainComplex\endcsname\trChainComplex
\renewcommand{\trChainComplex}{\BookTerm{T0122}{\csname BookOriginal@trChainComplex\endcsname}}
\expandafter\def\csname BookIndex@T0123\endcsname{lianyingshe@链映射（chain map）}
\expandafter\let\csname BookOriginal@trChainMap\endcsname\trChainMap
\renewcommand{\trChainMap}{\BookTerm{T0123}{\csname BookOriginal@trChainMap\endcsname}}
\expandafter\def\csname BookIndex@T0124\endcsname{bishanglian@闭上链（cocycle）}
\expandafter\let\csname BookOriginal@trCocycle\endcsname\trCocycle
\renewcommand{\trCocycle}{\BookTerm{T0124}{\csname BookOriginal@trCocycle\endcsname}}
\expandafter\def\csname BookIndex@T0125\endcsname{shanglianfuxing@上链复形（cochain complex）}
\expandafter\let\csname BookOriginal@trCochainComplex\endcsname\trCochainComplex
\renewcommand{\trCochainComplex}{\BookTerm{T0125}{\csname BookOriginal@trCochainComplex\endcsname}}
\expandafter\def\csname BookIndex@T0126\endcsname{bianyuansuanzi@边缘算子（boundary operator）}
\expandafter\let\csname BookOriginal@trBoundaryOperator\endcsname\trBoundaryOperator
\renewcommand{\trBoundaryOperator}{\BookTerm{T0126}{\csname BookOriginal@trBoundaryOperator\endcsname}}
\expandafter\def\csname BookIndex@T0127\endcsname{weifen@微分（differential）}
\expandafter\let\csname BookOriginal@trDifferential\endcsname\trDifferential
\renewcommand{\trDifferential}{\BookTerm{T0127}{\csname BookOriginal@trDifferential\endcsname}}
\expandafter\def\csname BookIndex@T0128\endcsname{beidishu@贝蒂数（betti number）}
\expandafter\let\csname BookOriginal@trBettiNumber\endcsname\trBettiNumber
\renewcommand{\trBettiNumber}{\BookTerm{T0128}{\csname BookOriginal@trBettiNumber\endcsname}}
\expandafter\def\csname BookIndex@T0129\endcsname{oulashixingshu@欧拉示性数（Euler characteristic）}
\expandafter\let\csname BookOriginal@trEulerCharacteristic\endcsname\trEulerCharacteristic
\renewcommand{\trEulerCharacteristic}{\BookTerm{T0129}{\csname BookOriginal@trEulerCharacteristic\endcsname}}
\expandafter\def\csname BookIndex@T0130\endcsname{pangjialaiduiou@庞加莱对偶（Poincaré duality）}
\expandafter\let\csname BookOriginal@trPoincareDuality\endcsname\trPoincareDuality
\renewcommand{\trPoincareDuality}{\BookTerm{T0130}{\csname BookOriginal@trPoincareDuality\endcsname}}
\expandafter\def\csname BookIndex@T0131\endcsname{kunnethgongshi@Künneth公式（Künneth formula）}
\expandafter\let\csname BookOriginal@trKunnethFormula\endcsname\trKunnethFormula
\renewcommand{\trKunnethFormula}{\BookTerm{T0131}{\csname BookOriginal@trKunnethFormula\endcsname}}
\expandafter\def\csname BookIndex@T0132\endcsname{barfenjie@bar~分解（bar resolution）}
\expandafter\let\csname BookOriginal@trBarResolution\endcsname\trBarResolution
\renewcommand{\trBarResolution}{\BookTerm{T0132}{\csname BookOriginal@trBarResolution\endcsname}}
\expandafter\def\csname BookIndex@T0133\endcsname{zhenghe@正合（exact）}
\expandafter\let\csname BookOriginal@trExact\endcsname\trExact
\renewcommand{\trExact}{\BookTerm{T0133}{\csname BookOriginal@trExact\endcsname}}
\expandafter\def\csname BookIndex@T0134\endcsname{zhangzhenghelie@长正合列（long exact sequence）}
\expandafter\let\csname BookOriginal@trLongExactSequence\endcsname\trLongExactSequence
\renewcommand{\trLongExactSequence}{\BookTerm{T0134}{\csname BookOriginal@trLongExactSequence\endcsname}}
\expandafter\def\csname BookIndex@T0135\endcsname{duanzhenghelie@短正合列（short exact sequence）}
\expandafter\let\csname BookOriginal@trShortExactSequence\endcsname\trShortExactSequence
\renewcommand{\trShortExactSequence}{\BookTerm{T0135}{\csname BookOriginal@trShortExactSequence\endcsname}}
\expandafter\def\csname BookIndex@T0136\endcsname{taishe@态射（morphism）}
\expandafter\let\csname BookOriginal@trMorphism\endcsname\trMorphism
\renewcommand{\trMorphism}{\BookTerm{T0136}{\csname BookOriginal@trMorphism\endcsname}}
\expandafter\def\csname BookIndex@T0137\endcsname{tongtai@同态（homomorphism）}
\expandafter\let\csname BookOriginal@trHomomorphism\endcsname\trHomomorphism
\renewcommand{\trHomomorphism}{\BookTerm{T0137}{\csname BookOriginal@trHomomorphism\endcsname}}
\expandafter\def\csname BookIndex@T0138\endcsname{dianfan@典范（canonical）}
\expandafter\let\csname BookOriginal@trCanonical\endcsname\trCanonical
\renewcommand{\trCanonical}{\BookTerm{T0138}{\csname BookOriginal@trCanonical\endcsname}}
\expandafter\def\csname BookIndex@T0139\endcsname{hanzi@函子（functor）}
\expandafter\let\csname BookOriginal@trFunctor\endcsname\trFunctor
\renewcommand{\trFunctor}{\BookTerm{T0139}{\csname BookOriginal@trFunctor\endcsname}}
\expandafter\def\csname BookIndex@T0140\endcsname{hanzixing@函子性（functoriality）}
\expandafter\let\csname BookOriginal@trFunctoriality\endcsname\trFunctoriality
\renewcommand{\trFunctoriality}{\BookTerm{T0140}{\csname BookOriginal@trFunctoriality\endcsname}}
\expandafter\def\csname BookIndex@T0141\endcsname{fanchou@范畴（category）}
\expandafter\let\csname BookOriginal@trCategory\endcsname\trCategory
\renewcommand{\trCategory}{\BookTerm{T0141}{\csname BookOriginal@trCategory\endcsname}}
\expandafter\def\csname BookIndex@T0142\endcsname{duixiang@对象（object）}
\expandafter\let\csname BookOriginal@trObject\endcsname\trObject
\renewcommand{\trObject}{\BookTerm{T0142}{\csname BookOriginal@trObject\endcsname}}
\expandafter\def\csname BookIndex@T0143\endcsname{peibian@配边（cobordism）}
\expandafter\let\csname BookOriginal@trCobordism\endcsname\trCobordism
\renewcommand{\trCobordism}{\BookTerm{T0143}{\csname BookOriginal@trCobordism\endcsname}}
\expandafter\def\csname BookIndex@T0144\endcsname{dingxiang@定向（orientation）}
\expandafter\let\csname BookOriginal@trOrientation\endcsname\trOrientation
\renewcommand{\trOrientation}{\BookTerm{T0144}{\csname BookOriginal@trOrientation\endcsname}}
\expandafter\def\csname BookIndex@T0145\endcsname{kedingxiang@可定向（orientable）}
\expandafter\let\csname BookOriginal@trOrientable\endcsname\trOrientable
\renewcommand{\trOrientable}{\BookTerm{T0145}{\csname BookOriginal@trOrientable\endcsname}}
\expandafter\def\csname BookIndex@T0146\endcsname{yudingxiang@余定向（coorientation）}
\expandafter\let\csname BookOriginal@trCoorientation\endcsname\trCoorientation
\renewcommand{\trCoorientation}{\BookTerm{T0146}{\csname BookOriginal@trCoorientation\endcsname}}
\expandafter\def\csname BookIndex@T0147\endcsname{fenci@分次（graded）}
\expandafter\let\csname BookOriginal@trGraded\endcsname\trGraded
\renewcommand{\trGraded}{\BookTerm{T0147}{\csname BookOriginal@trGraded\endcsname}}
\expandafter\def\csname BookIndex@T0148\endcsname{mo@模（module）}
\expandafter\let\csname BookOriginal@trModule\endcsname\trModule
\renewcommand{\trModule}{\BookTerm{T0148}{\csname BookOriginal@trModule\endcsname}}
\expandafter\def\csname BookIndex@T0149\endcsname{shuangmo@双模（bimodule）}
\expandafter\let\csname BookOriginal@trBimodule\endcsname\trBimodule
\renewcommand{\trBimodule}{\BookTerm{T0149}{\csname BookOriginal@trBimodule\endcsname}}
\expandafter\def\csname BookIndex@T0150\endcsname{fugai@覆盖（covering）}
\expandafter\let\csname BookOriginal@trCovering\endcsname\trCovering
\renewcommand{\trCovering}{\BookTerm{T0150}{\csname BookOriginal@trCovering\endcsname}}
\expandafter\def\csname BookIndex@T0151\endcsname{tonglun@同伦（homotopy）}
\expandafter\let\csname BookOriginal@trHomotopy\endcsname\trHomotopy
\renewcommand{\trHomotopy}{\BookTerm{T0151}{\csname BookOriginal@trHomotopy\endcsname}}
\expandafter\def\csname BookIndex@T0152\endcsname{xingbianshousuo@形变收缩（deformation retract）}
\expandafter\let\csname BookOriginal@trDeformationRetract\endcsname\trDeformationRetract
\renewcommand{\trDeformationRetract}{\BookTerm{T0152}{\csname BookOriginal@trDeformationRetract\endcsname}}
\expandafter\def\csname BookIndex@T0153\endcsname{chazhi@插值（interpolation）}
\expandafter\let\csname BookOriginal@trInterpolation\endcsname\trInterpolation
\renewcommand{\trInterpolation}{\BookTerm{T0153}{\csname BookOriginal@trInterpolation\endcsname}}
\expandafter\def\csname BookIndex@T0154\endcsname{xingbian@形变（deformation）}
\expandafter\let\csname BookOriginal@trDeformation\endcsname\trDeformation
\renewcommand{\trDeformation}{\BookTerm{T0154}{\csname BookOriginal@trDeformation\endcsname}}
\expandafter\def\csname BookIndex@T0155\endcsname{tonglun@同伦（homotopic）}
\expandafter\let\csname BookOriginal@trHomotopic\endcsname\trHomotopic
\renewcommand{\trHomotopic}{\BookTerm{T0155}{\csname BookOriginal@trHomotopic\endcsname}}
\expandafter\def\csname BookIndex@T0156\endcsname{tonglundengjia@同伦等价（homotopy equivalence）}
\expandafter\let\csname BookOriginal@trHomotopyEquivalence\endcsname\trHomotopyEquivalence
\renewcommand{\trHomotopyEquivalence}{\BookTerm{T0156}{\csname BookOriginal@trHomotopyEquivalence\endcsname}}
\expandafter\def\csname BookIndex@T0157\endcsname{liantonglun@链同伦（chain homotopy）}
\expandafter\let\csname BookOriginal@trChainHomotopy\endcsname\trChainHomotopy
\renewcommand{\trChainHomotopy}{\BookTerm{T0157}{\csname BookOriginal@trChainHomotopy\endcsname}}
\expandafter\def\csname BookIndex@T0158\endcsname{shousuotonglun@收缩同伦（contracting homotopy）}
\expandafter\let\csname BookOriginal@trContractingHomotopy\endcsname\trContractingHomotopy
\renewcommand{\trContractingHomotopy}{\BookTerm{T0158}{\csname BookOriginal@trContractingHomotopy\endcsname}}
\expandafter\def\csname BookIndex@T0159\endcsname{tonggou@同构（isomorphism）}
\expandafter\let\csname BookOriginal@trIsomorphism\endcsname\trIsomorphism
\renewcommand{\trIsomorphism}{\BookTerm{T0159}{\csname BookOriginal@trIsomorphism\endcsname}}
\expandafter\def\csname BookIndex@T0160\endcsname{zitongtai@自同态（endomorphism）}
\expandafter\let\csname BookOriginal@trEndomorphism\endcsname\trEndomorphism
\renewcommand{\trEndomorphism}{\BookTerm{T0160}{\csname BookOriginal@trEndomorphism\endcsname}}
\expandafter\def\csname BookIndex@T0161\endcsname{zitonggou@自同构（automorphism）}
\expandafter\let\csname BookOriginal@trAutomorphism\endcsname\trAutomorphism
\renewcommand{\trAutomorphism}{\BookTerm{T0161}{\csname BookOriginal@trAutomorphism\endcsname}}
\expandafter\def\csname BookIndex@T0162\endcsname{tonggou@同构（isomorphic）}
\expandafter\let\csname BookOriginal@trIsomorphic\endcsname\trIsomorphic
\renewcommand{\trIsomorphic}{\BookTerm{T0162}{\csname BookOriginal@trIsomorphic\endcsname}}
\expandafter\def\csname BookIndex@T0163\endcsname{he@核（kernel）}
\expandafter\let\csname BookOriginal@trKernel\endcsname\trKernel
\renewcommand{\trKernel}{\BookTerm{T0163}{\csname BookOriginal@trKernel\endcsname}}
\expandafter\def\csname BookIndex@T0164\endcsname{keni@可逆（invertible）}
\expandafter\let\csname BookOriginal@trInvertible\endcsname\trInvertible
\renewcommand{\trInvertible}{\BookTerm{T0164}{\csname BookOriginal@trInvertible\endcsname}}
\expandafter\def\csname BookIndex@T0165\endcsname{yuhe@余核（cokernel）}
\expandafter\let\csname BookOriginal@trCokernel\endcsname\trCokernel
\renewcommand{\trCokernel}{\BookTerm{T0165}{\csname BookOriginal@trCokernel\endcsname}}
\expandafter\def\csname BookIndex@T0166\endcsname{xiang@像（image）}
\expandafter\let\csname BookOriginal@trImage\endcsname\trImage
\renewcommand{\trImage}{\BookTerm{T0166}{\csname BookOriginal@trImage\endcsname}}
\expandafter\def\csname BookIndex@T0167\endcsname{yuanxiang@原像（preimage）}
\expandafter\let\csname BookOriginal@trPreimage\endcsname\trPreimage
\renewcommand{\trPreimage}{\BookTerm{T0167}{\csname BookOriginal@trPreimage\endcsname}}
\expandafter\def\csname BookIndex@T0168\endcsname{danshe@单射（injective）}
\expandafter\let\csname BookOriginal@trInjective\endcsname\trInjective
\renewcommand{\trInjective}{\BookTerm{T0168}{\csname BookOriginal@trInjective\endcsname}}
\expandafter\def\csname BookIndex@T0169\endcsname{manshe@满射（surjective）}
\expandafter\let\csname BookOriginal@trSurjective\endcsname\trSurjective
\renewcommand{\trSurjective}{\BookTerm{T0169}{\csname BookOriginal@trSurjective\endcsname}}
\expandafter\def\csname BookIndex@T0170\endcsname{shuangshe@双射（bijective）}
\expandafter\let\csname BookOriginal@trBijective\endcsname\trBijective
\renewcommand{\trBijective}{\BookTerm{T0170}{\csname BookOriginal@trBijective\endcsname}}
\expandafter\def\csname BookIndex@T0171\endcsname{zhangliangji@张量积（tensor product）}
\expandafter\let\csname BookOriginal@trTensorProduct\endcsname\trTensorProduct
\renewcommand{\trTensorProduct}{\BookTerm{T0171}{\csname BookOriginal@trTensorProduct\endcsname}}
\expandafter\def\csname BookIndex@T0172\endcsname{nao@挠（torsion）}
\expandafter\let\csname BookOriginal@trTorsion\endcsname\trTorsion
\renewcommand{\trTorsion}{\BookTerm{T0172}{\csname BookOriginal@trTorsion\endcsname}}
\expandafter\def\csname BookIndex@T0173\endcsname{zhihe@直和（direct sum）}
\expandafter\let\csname BookOriginal@trDirectSum\endcsname\trDirectSum
\renewcommand{\trDirectSum}{\BookTerm{T0173}{\csname BookOriginal@trDirectSum\endcsname}}
\expandafter\def\csname BookIndex@T0174\endcsname{zhihexiang@直和项（direct summand）}
\expandafter\let\csname BookOriginal@trDirectSummand\endcsname\trDirectSummand
\renewcommand{\trDirectSummand}{\BookTerm{T0174}{\csname BookOriginal@trDirectSummand\endcsname}}
\expandafter\def\csname BookIndex@T0175\endcsname{jubuxishu@局部系数（local coefficient）}
\expandafter\let\csname BookOriginal@trLocalCoefficient\endcsname\trLocalCoefficient
\renewcommand{\trLocalCoefficient}{\BookTerm{T0175}{\csname BookOriginal@trLocalCoefficient\endcsname}}
\expandafter\def\csname BookIndex@T0176\endcsname{xin@辛（symplectic）}
\expandafter\let\csname BookOriginal@trSymplectic\endcsname\trSymplectic
\renewcommand{\trSymplectic}{\BookTerm{T0176}{\csname BookOriginal@trSymplectic\endcsname}}
\expandafter\def\csname BookIndex@T0177\endcsname{xinxingshi@辛形式（symplectic form）}
\expandafter\let\csname BookOriginal@trSymplecticForm\endcsname\trSymplecticForm
\renewcommand{\trSymplecticForm}{\BookTerm{T0177}{\csname BookOriginal@trSymplecticForm\endcsname}}
\expandafter\def\csname BookIndex@T0178\endcsname{fanduicheng@反对称（antisymmetric）}
\expandafter\let\csname BookOriginal@trAntisymmetric\endcsname\trAntisymmetric
\renewcommand{\trAntisymmetric}{\BookTerm{T0178}{\csname BookOriginal@trAntisymmetric\endcsname}}
\expandafter\def\csname BookIndex@T0179\endcsname{aiermite@埃尔米特（hermitian）}
\expandafter\let\csname BookOriginal@trHermitian\endcsname\trHermitian
\renewcommand{\trHermitian}{\BookTerm{T0179}{\csname BookOriginal@trHermitian\endcsname}}
\expandafter\def\csname BookIndex@T0180\endcsname{kailei@凯勒（Kähler）}
\expandafter\let\csname BookOriginal@trKaehler\endcsname\trKaehler
\renewcommand{\trKaehler}{\BookTerm{T0180}{\csname BookOriginal@trKaehler\endcsname}}
\expandafter\def\csname BookIndex@T0181\endcsname{dengju@等距（isometry）}
\expandafter\let\csname BookOriginal@trIsometry\endcsname\trIsometry
\renewcommand{\trIsometry}{\BookTerm{T0181}{\csname BookOriginal@trIsometry\endcsname}}
\expandafter\def\csname BookIndex@T0182\endcsname{tezhengzhi@特征值（eigenvalue）}
\expandafter\let\csname BookOriginal@trEigenvalue\endcsname\trEigenvalue
\renewcommand{\trEigenvalue}{\BookTerm{T0182}{\csname BookOriginal@trEigenvalue\endcsname}}
\expandafter\def\csname BookIndex@T0183\endcsname{tezhengbiao@特征标（character）}
\expandafter\let\csname BookOriginal@trCharacter\endcsname\trCharacter
\renewcommand{\trCharacter}{\BookTerm{T0183}{\csname BookOriginal@trCharacter\endcsname}}
\expandafter\def\csname BookIndex@T0184\endcsname{zhongshu@重数（multiplicity）}
\expandafter\let\csname BookOriginal@trMultiplicity\endcsname\trMultiplicity
\renewcommand{\trMultiplicity}{\BookTerm{T0184}{\csname BookOriginal@trMultiplicity\endcsname}}
\expandafter\def\csname BookIndex@T0185\endcsname{tezhengzikongjian@特征子空间（eigenspace）}
\expandafter\let\csname BookOriginal@trEigenspace\endcsname\trEigenspace
\renewcommand{\trEigenspace}{\BookTerm{T0185}{\csname BookOriginal@trEigenspace\endcsname}}
\expandafter\def\csname BookIndex@T0186\endcsname{keduijiaohua@可对角化（diagonalizable）}
\expandafter\let\csname BookOriginal@trDiagonalizable\endcsname\trDiagonalizable
\renewcommand{\trDiagonalizable}{\BookTerm{T0186}{\csname BookOriginal@trDiagonalizable\endcsname}}
\expandafter\def\csname BookIndex@T0187\endcsname{fuhua@复化（complexification）}
\expandafter\let\csname BookOriginal@trComplexification\endcsname\trComplexification
\renewcommand{\trComplexification}{\BookTerm{T0187}{\csname BookOriginal@trComplexification\endcsname}}
\expandafter\def\csname BookIndex@T0188\endcsname{danzhiyingshe@单值映射（monodromy）}
\expandafter\let\csname BookOriginal@trMonodromy\endcsname\trMonodromy
\renewcommand{\trMonodromy}{\BookTerm{T0188}{\csname BookOriginal@trMonodromy\endcsname}}
\expandafter\def\csname BookIndex@T0189\endcsname{bi@闭（closed）}
\expandafter\let\csname BookOriginal@trClosed\endcsname\trClosed
\renewcommand{\trClosed}{\BookTerm{T0189}{\csname BookOriginal@trClosed\endcsname}}
\expandafter\def\csname BookIndex@T0190\endcsname{xingshi@形式（form）}
\expandafter\let\csname BookOriginal@trForm\endcsname\trForm
\renewcommand{\trForm}{\BookTerm{T0190}{\csname BookOriginal@trForm\endcsname}}
\expandafter\def\csname BookIndex@T0191\endcsname{tiji@体积（volume）}
\expandafter\let\csname BookOriginal@trVolume\endcsname\trVolume
\renewcommand{\trVolume}{\BookTerm{T0191}{\csname BookOriginal@trVolume\endcsname}}
\expandafter\def\csname BookIndex@T0192\endcsname{jinfu@近复（almost complex）}
\expandafter\let\csname BookOriginal@trAlmostComplex\endcsname\trAlmostComplex
\renewcommand{\trAlmostComplex}{\BookTerm{T0192}{\csname BookOriginal@trAlmostComplex\endcsname}}
\expandafter\def\csname BookIndex@T0193\endcsname{fujiegou@复结构（complex structure）}
\expandafter\let\csname BookOriginal@trComplexStructure\endcsname\trComplexStructure
\renewcommand{\trComplexStructure}{\BookTerm{T0193}{\csname BookOriginal@trComplexStructure\endcsname}}
\expandafter\def\csname BookIndex@T0194\endcsname{xiangrong@相容（compatible）}
\expandafter\let\csname BookOriginal@trCompatible\endcsname\trCompatible
\renewcommand{\trCompatible}{\BookTerm{T0194}{\csname BookOriginal@trCompatible\endcsname}}
\expandafter\def\csname BookIndex@T0195\endcsname{hamidun@哈密顿（Hamiltonian）}
\expandafter\let\csname BookOriginal@trHamiltonian\endcsname\trHamiltonian
\renewcommand{\trHamiltonian}{\BookTerm{T0195}{\csname BookOriginal@trHamiltonian\endcsname}}
\expandafter\def\csname BookIndex@T0196\endcsname{hamidunxiangliangchang@哈密顿向量场（Hamiltonian vector field）}
\expandafter\let\csname BookOriginal@trHamiltonianVectorField\endcsname\trHamiltonianVectorField
\renewcommand{\trHamiltonianVectorField}{\BookTerm{T0196}{\csname BookOriginal@trHamiltonianVectorField\endcsname}}
\expandafter\def\csname BookIndex@T0197\endcsname{hamiduntonghen@哈密顿同痕（Hamiltonian isotopy）}
\expandafter\let\csname BookOriginal@trHamiltonianIsotopy\endcsname\trHamiltonianIsotopy
\renewcommand{\trHamiltonianIsotopy}{\BookTerm{T0197}{\csname BookOriginal@trHamiltonianIsotopy\endcsname}}
\expandafter\def\csname BookIndex@T0198\endcsname{zuoyongfanhan@作用泛函（action functional）}
\expandafter\let\csname BookOriginal@trActionFunctional\endcsname\trActionFunctional
\renewcommand{\trActionFunctional}{\BookTerm{T0198}{\csname BookOriginal@trActionFunctional\endcsname}}
\expandafter\def\csname BookIndex@T0199\endcsname{fuluoertongdiao@弗洛尔同调（Floer homology）}
\expandafter\let\csname BookOriginal@trFloerHomology\endcsname\trFloerHomology
\renewcommand{\trFloerHomology}{\BookTerm{T0199}{\csname BookOriginal@trFloerHomology\endcsname}}
\expandafter\def\csname BookIndex@T0200\endcsname{fuluoershangtongdiao@弗洛尔上同调（Floer cohomology）}
\expandafter\let\csname BookOriginal@trFloerCohomology\endcsname\trFloerCohomology
\renewcommand{\trFloerCohomology}{\BookTerm{T0200}{\csname BookOriginal@trFloerCohomology\endcsname}}
\expandafter\def\csname BookIndex@T0201\endcsname{fuluoerfuxing@弗洛尔复形（Floer complex）}
\expandafter\let\csname BookOriginal@trFloerComplex\endcsname\trFloerComplex
\renewcommand{\trFloerComplex}{\BookTerm{T0201}{\csname BookOriginal@trFloerComplex\endcsname}}
\expandafter\def\csname BookIndex@T0202\endcsname{fuluoershanglianqun@弗洛尔上链群（Floer cochain group）}
\expandafter\let\csname BookOriginal@trFloerCochainGroup\endcsname\trFloerCochainGroup
\renewcommand{\trFloerCochainGroup}{\BookTerm{T0202}{\csname BookOriginal@trFloerCochainGroup\endcsname}}
\expandafter\def\csname BookIndex@T0203\endcsname{fuluoerfangcheng@弗洛尔方程（Floer equation）}
\expandafter\let\csname BookOriginal@trFloerEquation\endcsname\trFloerEquation
\renewcommand{\trFloerEquation}{\BookTerm{T0203}{\csname BookOriginal@trFloerEquation\endcsname}}
\expandafter\def\csname BookIndex@T0204\endcsname{masiluofuzhibiao@马斯洛夫指标（Maslov index）}
\expandafter\let\csname BookOriginal@trMaslovIndex\endcsname\trMaslovIndex
\renewcommand{\trMaslovIndex}{\BookTerm{T0204}{\csname BookOriginal@trMaslovIndex\endcsname}}
\expandafter\def\csname BookIndex@T0205\endcsname{masiluofulei@马斯洛夫类（Maslov class）}
\expandafter\let\csname BookOriginal@trMaslovClass\endcsname\trMaslovClass
\renewcommand{\trMaslovClass}{\BookTerm{T0205}{\csname BookOriginal@trMaslovClass\endcsname}}
\expandafter\def\csname BookIndex@T0206\endcsname{bilian@闭链（cycle）}
\expandafter\let\csname BookOriginal@trCycle\endcsname\trCycle
\renewcommand{\trCycle}{\BookTerm{T0206}{\csname BookOriginal@trCycle\endcsname}}
\expandafter\def\csname BookIndex@T0207\endcsname{huanlu@环路（loop）}
\expandafter\let\csname BookOriginal@trLoop\endcsname\trLoop
\renewcommand{\trLoop}{\BookTerm{T0207}{\csname BookOriginal@trLoop\endcsname}}
\expandafter\def\csname BookIndex@T0208\endcsname{huanlukongjian@环路空间（loop space）}
\expandafter\let\csname BookOriginal@trLoopSpace\endcsname\trLoopSpace
\renewcommand{\trLoopSpace}{\BookTerm{T0208}{\csname BookOriginal@trLoopSpace\endcsname}}
\expandafter\def\csname BookIndex@T0209\endcsname{kesuo@可缩（contractible）}
\expandafter\let\csname BookOriginal@trContractible\endcsname\trContractible
\renewcommand{\trContractible}{\BookTerm{T0209}{\csname BookOriginal@trContractible\endcsname}}
\expandafter\def\csname BookIndex@T0210\endcsname{zhouqi@周期（periodic）}
\expandafter\let\csname BookOriginal@trPeriodic\endcsname\trPeriodic
\renewcommand{\trPeriodic}{\BookTerm{T0210}{\csname BookOriginal@trPeriodic\endcsname}}
\expandafter\def\csname BookIndex@T0211\endcsname{budongdian@不动点（fixed point）}
\expandafter\let\csname BookOriginal@trFixedPoint\endcsname\trFixedPoint
\renewcommand{\trFixedPoint}{\BookTerm{T0211}{\csname BookOriginal@trFixedPoint\endcsname}}
\expandafter\def\csname BookIndex@T0212\endcsname{chenlei@陈类（Chern class）}
\expandafter\let\csname BookOriginal@trChernClass\endcsname\trChernClass
\renewcommand{\trChernClass}{\BookTerm{T0212}{\csname BookOriginal@trChernClass\endcsname}}
\expandafter\def\csname BookIndex@T0213\endcsname{diyichenlei@第一陈类（first Chern class）}
\expandafter\let\csname BookOriginal@trFirstChernClass\endcsname\trFirstChernClass
\renewcommand{\trFirstChernClass}{\BookTerm{T0213}{\csname BookOriginal@trFirstChernClass\endcsname}}
\expandafter\def\csname BookIndex@T0214\endcsname{zhongyan@重言（tautological）}
\expandafter\let\csname BookOriginal@trTautological\endcsname\trTautological
\renewcommand{\trTautological}{\BookTerm{T0214}{\csname BookOriginal@trTautological\endcsname}}
\expandafter\def\csname BookIndex@T0215\endcsname{lagelangrigelasimanliuxing@拉格朗日--格拉斯曼流形（Lagrangian grassmannian）}
\expandafter\let\csname BookOriginal@trLagrangianGrassmannian\endcsname\trLagrangianGrassmannian
\renewcommand{\trLagrangianGrassmannian}{\BookTerm{T0215}{\csname BookOriginal@trLagrangianGrassmannian\endcsname}}
\expandafter\def\csname BookIndex@T0216\endcsname{leirangdeshi@勒让德式（legendrian）}
\expandafter\let\csname BookOriginal@trLegendrian\endcsname\trLegendrian
\renewcommand{\trLegendrian}{\BookTerm{T0216}{\csname BookOriginal@trLegendrian\endcsname}}
\expandafter\def\csname BookIndex@T0217\endcsname{zhengjiao@正交（orthogonal）}
\expandafter\let\csname BookOriginal@trOrthogonal\endcsname\trOrthogonal
\renewcommand{\trOrthogonal}{\BookTerm{T0217}{\csname BookOriginal@trOrthogonal\endcsname}}
\expandafter\def\csname BookIndex@T0218\endcsname{biaozhunzhengjiao@标准正交（orthonormal）}
\expandafter\let\csname BookOriginal@trOrthonormal\endcsname\trOrthonormal
\renewcommand{\trOrthonormal}{\BookTerm{T0218}{\csname BookOriginal@trOrthonormal\endcsname}}
\expandafter\def\csname BookIndex@T0219\endcsname{neiji@内积（inner product）}
\expandafter\let\csname BookOriginal@trInnerProduct\endcsname\trInnerProduct
\renewcommand{\trInnerProduct}{\BookTerm{T0219}{\csname BookOriginal@trInnerProduct\endcsname}}
\expandafter\def\csname BookIndex@T0220\endcsname{puliu@谱流（spectral flow）}
\expandafter\let\csname BookOriginal@trSpectralFlow\endcsname\trSpectralFlow
\renewcommand{\trSpectralFlow}{\BookTerm{T0220}{\csname BookOriginal@trSpectralFlow\endcsname}}
\expandafter\def\csname BookIndex@T0221\endcsname{chuanyuedian@穿越点（crossing）}
\expandafter\let\csname BookOriginal@trCrossing\endcsname\trCrossing
\renewcommand{\trCrossing}{\BookTerm{T0221}{\csname BookOriginal@trCrossing\endcsname}}
\expandafter\def\csname BookIndex@T0222\endcsname{weiquanchunquxian@伪全纯曲线（pseudoholomorphic curve）}
\expandafter\let\csname BookOriginal@trPseudoholomorphicCurve\endcsname\trPseudoholomorphicCurve
\renewcommand{\trPseudoholomorphicCurve}{\BookTerm{T0222}{\csname BookOriginal@trPseudoholomorphicCurve\endcsname}}
\expandafter\def\csname BookIndex@T0223\endcsname{quanchun@全纯（holomorphic）}
\expandafter\let\csname BookOriginal@trHolomorphic\endcsname\trHolomorphic
\renewcommand{\trHolomorphic}{\BookTerm{T0223}{\csname BookOriginal@trHolomorphic\endcsname}}
\expandafter\def\csname BookIndex@T0224\endcsname{weiquanchun@伪全纯（pseudoholomorphic）}
\expandafter\let\csname BookOriginal@trPseudoholomorphic\endcsname\trPseudoholomorphic
\renewcommand{\trPseudoholomorphic}{\BookTerm{T0224}{\csname BookOriginal@trPseudoholomorphic\endcsname}}
\expandafter\def\csname BookIndex@T0225\endcsname{kexiliman@柯西--黎曼（Cauchy–Riemann）}
\expandafter\let\csname BookOriginal@trCauchyRiemann\endcsname\trCauchyRiemann
\renewcommand{\trCauchyRiemann}{\BookTerm{T0225}{\csname BookOriginal@trCauchyRiemann\endcsname}}
\expandafter\def\csname BookIndex@T0226\endcsname{nengliang@能量（energy）}
\expandafter\let\csname BookOriginal@trEnergy\endcsname\trEnergy
\renewcommand{\trEnergy}{\BookTerm{T0226}{\csname BookOriginal@trEnergy\endcsname}}
\expandafter\def\csname BookIndex@T0227\endcsname{yanxuyingshe@延续映射（continuation map）}
\expandafter\let\csname BookOriginal@trContinuationMap\endcsname\trContinuationMap
\renewcommand{\trContinuationMap}{\BookTerm{T0227}{\csname BookOriginal@trContinuationMap\endcsname}}
\expandafter\def\csname BookIndex@T0228\endcsname{raodong@扰动（perturbation）}
\expandafter\let\csname BookOriginal@trPerturbation\endcsname\trPerturbation
\renewcommand{\trPerturbation}{\BookTerm{T0228}{\csname BookOriginal@trPerturbation\endcsname}}
\expandafter\def\csname BookIndex@T0229\endcsname{xianxinghua@线性化（linearization）}
\expandafter\let\csname BookOriginal@trLinearization\endcsname\trLinearization
\renewcommand{\trLinearization}{\BookTerm{T0229}{\csname BookOriginal@trLinearization\endcsname}}
\expandafter\def\csname BookIndex@T0230\endcsname{banaheliuxing@巴拿赫流形（Banach manifold）}
\expandafter\let\csname BookOriginal@trBanachManifold\endcsname\trBanachManifold
\renewcommand{\trBanachManifold}{\BookTerm{T0230}{\csname BookOriginal@trBanachManifold\endcsname}}
\expandafter\def\csname BookIndex@T0231\endcsname{jingquexin@精确辛（exact symplectic）}
\expandafter\let\csname BookOriginal@trExactSymplectic\endcsname\trExactSymplectic
\renewcommand{\trExactSymplectic}{\BookTerm{T0231}{\csname BookOriginal@trExactSymplectic\endcsname}}
\expandafter\def\csname BookIndex@T0232\endcsname{tu@凸（convex）}
\expandafter\let\csname BookOriginal@trConvex\endcsname\trConvex
\renewcommand{\trConvex}{\BookTerm{T0232}{\csname BookOriginal@trConvex\endcsname}}
\expandafter\def\csname BookIndex@T0233\endcsname{jiechuxingbianjie@接触型边界（contact type boundary）}
\expandafter\let\csname BookOriginal@trContactTypeBoundary\endcsname\trContactTypeBoundary
\renewcommand{\trContactTypeBoundary}{\BookTerm{T0233}{\csname BookOriginal@trContactTypeBoundary\endcsname}}
\expandafter\def\csname BookIndex@T0234\endcsname{gongxingxin@共形辛（conformally symplectic）}
\expandafter\let\csname BookOriginal@trConformallySymplectic\endcsname\trConformallySymplectic
\renewcommand{\trConformallySymplectic}{\BookTerm{T0234}{\csname BookOriginal@trConformallySymplectic\endcsname}}
\expandafter\def\csname BookIndex@T0235\endcsname{jidazhiyuanli@极大值原理（maximum principle）}
\expandafter\let\csname BookOriginal@trMaximumPrinciple\endcsname\trMaximumPrinciple
\renewcommand{\trMaximumPrinciple}{\BookTerm{T0235}{\csname BookOriginal@trMaximumPrinciple\endcsname}}
\expandafter\def\csname BookIndex@T0236\endcsname{shangbanpingmian@上半平面（upper half plane）}
\expandafter\let\csname BookOriginal@trUpperHalfPlane\endcsname\trUpperHalfPlane
\renewcommand{\trUpperHalfPlane}{\BookTerm{T0236}{\csname BookOriginal@trUpperHalfPlane\endcsname}}
\expandafter\def\csname BookIndex@T0237\endcsname{shuangquanchun@双全纯（biholomorphic）}
\expandafter\let\csname BookOriginal@trBiholomorphic\endcsname\trBiholomorphic
\renewcommand{\trBiholomorphic}{\BookTerm{T0237}{\csname BookOriginal@trBiholomorphic\endcsname}}
\expandafter\def\csname BookIndex@T0238\endcsname{daibiaojibianjiedian@带标记边界点（pointed boundary）}
\expandafter\let\csname BookOriginal@trPointedBoundary\endcsname\trPointedBoundary
\renewcommand{\trPointedBoundary}{\BookTerm{T0238}{\csname BookOriginal@trPointedBoundary\endcsname}}
\expandafter\def\csname BookIndex@T0239\endcsname{rudian@入点（incoming point）}
\expandafter\let\csname BookOriginal@trIncomingPoint\endcsname\trIncomingPoint
\renewcommand{\trIncomingPoint}{\BookTerm{T0239}{\csname BookOriginal@trIncomingPoint\endcsname}}
\expandafter\def\csname BookIndex@T0240\endcsname{chudian@出点（outgoing point）}
\expandafter\let\csname BookOriginal@trOutgoingPoint\endcsname\trOutgoingPoint
\renewcommand{\trOutgoingPoint}{\BookTerm{T0240}{\csname BookOriginal@trOutgoingPoint\endcsname}}
\expandafter\def\csname BookIndex@T0241\endcsname{biaoqian@标签（label）}
\expandafter\let\csname BookOriginal@trLabel\endcsname\trLabel
\renewcommand{\trLabel}{\BookTerm{T0241}{\csname BookOriginal@trLabel\endcsname}}
\expandafter\def\csname BookIndex@T0242\endcsname{dai@带（strip）}
\expandafter\let\csname BookOriginal@trStrip\endcsname\trStrip
\renewcommand{\trStrip}{\BookTerm{T0242}{\csname BookOriginal@trStrip\endcsname}}
\expandafter\def\csname BookIndex@T0243\endcsname{yuanhuajiao@圆化角（round off corners）}
\expandafter\let\csname BookOriginal@trRoundOffCorners\endcsname\trRoundOffCorners
\renewcommand{\trRoundOffCorners}{\BookTerm{T0243}{\csname BookOriginal@trRoundOffCorners\endcsname}}
\expandafter\def\csname BookIndex@T0244\endcsname{fuleidehuomusuanzi@弗雷德霍姆算子（Fredholm operator）}
\expandafter\let\csname BookOriginal@trFredholmOperator\endcsname\trFredholmOperator
\renewcommand{\trFredholmOperator}{\BookTerm{T0244}{\csname BookOriginal@trFredholmOperator\endcsname}}
\expandafter\def\csname BookIndex@T0245\endcsname{banahekongjian@巴拿赫空间（Banach space）}
\expandafter\let\csname BookOriginal@trBanachSpace\endcsname\trBanachSpace
\renewcommand{\trBanachSpace}{\BookTerm{T0245}{\csname BookOriginal@trBanachSpace\endcsname}}
\expandafter\def\csname BookIndex@T0246\endcsname{xierbotekongjian@希尔伯特空间（Hilbert space）}
\expandafter\let\csname BookOriginal@trHilbertSpace\endcsname\trHilbertSpace
\renewcommand{\trHilbertSpace}{\BookTerm{T0246}{\csname BookOriginal@trHilbertSpace\endcsname}}
\expandafter\def\csname BookIndex@T0247\endcsname{suoboliefukongjian@索伯列夫空间（Sobolev space）}
\expandafter\let\csname BookOriginal@trSobolevSpace\endcsname\trSobolevSpace
\renewcommand{\trSobolevSpace}{\BookTerm{T0247}{\csname BookOriginal@trSobolevSpace\endcsname}}
\expandafter\def\csname BookIndex@T0248\endcsname{ruoshoulian@弱收敛（weak convergence）}
\expandafter\let\csname BookOriginal@trWeakConvergence\endcsname\trWeakConvergence
\renewcommand{\trWeakConvergence}{\BookTerm{T0248}{\csname BookOriginal@trWeakConvergence\endcsname}}
\expandafter\def\csname BookIndex@T0249\endcsname{shoulian@收敛（convergence）}
\expandafter\let\csname BookOriginal@trConvergence\endcsname\trConvergence
\renewcommand{\trConvergence}{\BookTerm{T0249}{\csname BookOriginal@trConvergence\endcsname}}
\expandafter\def\csname BookIndex@T0250\endcsname{tuoyuan@椭圆（elliptic）}
\expandafter\let\csname BookOriginal@trElliptic\endcsname\trElliptic
\renewcommand{\trElliptic}{\BookTerm{T0250}{\csname BookOriginal@trElliptic\endcsname}}
\expandafter\def\csname BookIndex@T0251\endcsname{tuoyuanzhengzexing@椭圆正则性（elliptic regularity）}
\expandafter\let\csname BookOriginal@trEllipticRegularity\endcsname\trEllipticRegularity
\renewcommand{\trEllipticRegularity}{\BookTerm{T0251}{\csname BookOriginal@trEllipticRegularity\endcsname}}
\expandafter\def\csname BookIndex@T0252\endcsname{youjie@有界（bounded）}
\expandafter\let\csname BookOriginal@trBounded\endcsname\trBounded
\renewcommand{\trBounded}{\BookTerm{T0252}{\csname BookOriginal@trBounded\endcsname}}
\expandafter\def\csname BookIndex@T0253\endcsname{yizhi@一致（uniform）}
\expandafter\let\csname BookOriginal@trUniform\endcsname\trUniform
\renewcommand{\trUniform}{\BookTerm{T0253}{\csname BookOriginal@trUniform\endcsname}}
\expandafter\def\csname BookIndex@T0254\endcsname{yinhanshudingli@隐函数定理（implicit function theorem）}
\expandafter\let\csname BookOriginal@trImplicitFunctionTheorem\endcsname\trImplicitFunctionTheorem
\renewcommand{\trImplicitFunctionTheorem}{\BookTerm{T0254}{\csname BookOriginal@trImplicitFunctionTheorem\endcsname}}
\expandafter\def\csname BookIndex@T0255\endcsname{zhengzexing@正则性（regularity）}
\expandafter\let\csname BookOriginal@trRegularity\endcsname\trRegularity
\renewcommand{\trRegularity}{\BookTerm{T0255}{\csname BookOriginal@trRegularity\endcsname}}
\expandafter\def\csname BookIndex@T0256\endcsname{xianyan@先验（apriori）}
\expandafter\let\csname BookOriginal@trApriori\endcsname\trApriori
\renewcommand{\trApriori}{\BookTerm{T0256}{\csname BookOriginal@trApriori\endcsname}}
\expandafter\def\csname BookIndex@T0257\endcsname{ainftyfanchou@$A_\infty$-范畴（\ensuremath{A_\infty}-category）}
\expandafter\let\csname BookOriginal@trACategory\endcsname\trACategory
\renewcommand{\trACategory}{\BookTerm{T0257}{\csname BookOriginal@trACategory\endcsname}}
\expandafter\def\csname BookIndex@T0258\endcsname{ainftyhanzi@$A_\infty$-函子（\ensuremath{A_\infty}-functor）}
\expandafter\let\csname BookOriginal@trAFunctor\endcsname\trAFunctor
\renewcommand{\trAFunctor}{\BookTerm{T0258}{\csname BookOriginal@trAFunctor\endcsname}}
\expandafter\def\csname BookIndex@T0259\endcsname{liwaixulie@例外序列（exceptional collection）}
\expandafter\let\csname BookOriginal@trExceptionalCollection\endcsname\trExceptionalCollection
\renewcommand{\trExceptionalCollection}{\BookTerm{T0259}{\csname BookOriginal@trExceptionalCollection\endcsname}}
\expandafter\def\csname BookIndex@T0260\endcsname{bianyi@变异（mutation）}
\expandafter\let\csname BookOriginal@trMutation\endcsname\trMutation
\renewcommand{\trMutation}{\BookTerm{T0260}{\csname BookOriginal@trMutation\endcsname}}
\expandafter\def\csname BookIndex@T0261\endcsname{qiumianduixiang@球面对象（spherical object）}
\expandafter\let\csname BookOriginal@trSphericalObject\endcsname\trSphericalObject
\renewcommand{\trSphericalObject}{\BookTerm{T0261}{\csname BookOriginal@trSphericalObject\endcsname}}
\expandafter\def\csname BookIndex@T0262\endcsname{luguo@滤过（filtration）}
\expandafter\let\csname BookOriginal@trFiltration\endcsname\trFiltration
\renewcommand{\trFiltration}{\BookTerm{T0262}{\csname BookOriginal@trFiltration\endcsname}}
\expandafter\def\csname BookIndex@T0263\endcsname{jiegouyingshe@结构映射（structure map）}
\expandafter\let\csname BookOriginal@trStructureMap\endcsname\trStructureMap
\renewcommand{\trStructureMap}{\BookTerm{T0263}{\csname BookOriginal@trStructureMap\endcsname}}
\expandafter\def\csname BookIndex@T0264\endcsname{fenchai@分拆（partition）}
\expandafter\let\csname BookOriginal@trPartition\endcsname\trPartition
\renewcommand{\trPartition}{\BookTerm{T0264}{\csname BookOriginal@trPartition\endcsname}}
\expandafter\def\csname BookIndex@T0265\endcsname{shangbianyuan@上边缘（coboundary）}
\expandafter\let\csname BookOriginal@trCoboundary\endcsname\trCoboundary
\renewcommand{\trCoboundary}{\BookTerm{T0265}{\csname BookOriginal@trCoboundary\endcsname}}
\expandafter\def\csname BookIndex@T0266\endcsname{jingxiangduicheng@镜像对称（mirror symmetry）}
\expandafter\let\csname BookOriginal@trMirrorSymmetry\endcsname\trMirrorSymmetry
\renewcommand{\trMirrorSymmetry}{\BookTerm{T0266}{\csname BookOriginal@trMirrorSymmetry\endcsname}}
\expandafter\def\csname BookIndex@T0267\endcsname{feiyao@非幺（nonunital）}
\expandafter\let\csname BookOriginal@trNonunital\endcsname\trNonunital
\renewcommand{\trNonunital}{\BookTerm{T0267}{\csname BookOriginal@trNonunital\endcsname}}
\expandafter\def\csname BookIndex@T0268\endcsname{yao@幺（unital）}
\expandafter\let\csname BookOriginal@trUnital\endcsname\trUnital
\renewcommand{\trUnital}{\BookTerm{T0268}{\csname BookOriginal@trUnital\endcsname}}
\expandafter\def\csname BookIndex@T0269\endcsname{yangeyao@严格幺（strictly unital）}
\expandafter\let\csname BookOriginal@trStrictlyUnital\endcsname\trStrictlyUnital
\renewcommand{\trStrictlyUnital}{\BookTerm{T0269}{\csname BookOriginal@trStrictlyUnital\endcsname}}
\expandafter\def\csname BookIndex@T0270\endcsname{shangtongdiaoyao@上同调幺（cunital）}
\expandafter\let\csname BookOriginal@trCunital\endcsname\trCunital
\renewcommand{\trCunital}{\BookTerm{T0270}{\csname BookOriginal@trCunital\endcsname}}
\expandafter\def\csname BookIndex@T0271\endcsname{cyouxian@c-有限（cfinite）}
\expandafter\let\csname BookOriginal@trCfinite\endcsname\trCfinite
\renewcommand{\trCfinite}{\BookTerm{T0271}{\csname BookOriginal@trCfinite\endcsname}}
\expandafter\def\csname BookIndex@T0272\endcsname{hengdengtaishe@恒等态射（identity morphism）}
\expandafter\let\csname BookOriginal@trIdentityMorphism\endcsname\trIdentityMorphism
\renewcommand{\trIdentityMorphism}{\BookTerm{T0272}{\csname BookOriginal@trIdentityMorphism\endcsname}}
\expandafter\def\csname BookIndex@T0273\endcsname{mideng@幂等（idempotent）}
\expandafter\let\csname BookOriginal@trIdempotent\endcsname\trIdempotent
\renewcommand{\trIdempotent}{\BookTerm{T0273}{\csname BookOriginal@trIdempotent\endcsname}}
\expandafter\def\csname BookIndex@T0274\endcsname{nitonggou@拟同构（quasiisomorphism）}
\expandafter\let\csname BookOriginal@trQuasiisomorphism\endcsname\trQuasiisomorphism
\renewcommand{\trQuasiisomorphism}{\BookTerm{T0274}{\csname BookOriginal@trQuasiisomorphism\endcsname}}
\expandafter\def\csname BookIndex@T0275\endcsname{nitonggou@拟同构（quasiisomorphic）}
\expandafter\let\csname BookOriginal@trQuasiisomorphic\endcsname\trQuasiisomorphic
\renewcommand{\trQuasiisomorphic}{\BookTerm{T0275}{\csname BookOriginal@trQuasiisomorphic\endcsname}}
\expandafter\def\csname BookIndex@T0276\endcsname{nibiaoshi@拟表示（quasi represent）}
\expandafter\let\csname BookOriginal@trQuasiRepresent\endcsname\trQuasiRepresent
\renewcommand{\trQuasiRepresent}{\BookTerm{T0276}{\csname BookOriginal@trQuasiRepresent\endcsname}}
\expandafter\def\csname BookIndex@T0277\endcsname{nidengjia@拟等价（quasiequivalence）}
\expandafter\let\csname BookOriginal@trQuasiequivalence\endcsname\trQuasiequivalence
\renewcommand{\trQuasiequivalence}{\BookTerm{T0277}{\csname BookOriginal@trQuasiequivalence\endcsname}}
\expandafter\def\csname BookIndex@T0278\endcsname{zhenghesanjiaoxing@正合三角形（exact triangle）}
\expandafter\let\csname BookOriginal@trExactTriangle\endcsname\trExactTriangle
\renewcommand{\trExactTriangle}{\BookTerm{T0278}{\csname BookOriginal@trExactTriangle\endcsname}}
\expandafter\def\csname BookIndex@T0279\endcsname{pingyi@平移（translation）}
\expandafter\let\csname BookOriginal@trTranslation\endcsname\trTranslation
\renewcommand{\trTranslation}{\BookTerm{T0279}{\csname BookOriginal@trTranslation\endcsname}}
\expandafter\def\csname BookIndex@T0280\endcsname{teyisanjiaoxing@特异三角形（distinguished triangle）}
\expandafter\let\csname BookOriginal@trDistinguishedTriangle\endcsname\trDistinguishedTriangle
\renewcommand{\trDistinguishedTriangle}{\BookTerm{T0280}{\csname BookOriginal@trDistinguishedTriangle\endcsname}}
\expandafter\def\csname BookIndex@T0281\endcsname{baoluo@包络（envelope）}
\expandafter\let\csname BookOriginal@trEnvelope\endcsname\trEnvelope
\renewcommand{\trEnvelope}{\BookTerm{T0281}{\csname BookOriginal@trEnvelope\endcsname}}
\expandafter\def\csname BookIndex@T0282\endcsname{sanjiaohua@三角化（triangulated）}
\expandafter\let\csname BookOriginal@trTriangulated\endcsname\trTriangulated
\renewcommand{\trTriangulated}{\BookTerm{T0282}{\csname BookOriginal@trTriangulated\endcsname}}
\expandafter\def\csname BookIndex@T0283\endcsname{niuqufuxing@扭曲复形（twisted complex）}
\expandafter\let\csname BookOriginal@trTwistedComplex\endcsname\trTwistedComplex
\renewcommand{\trTwistedComplex}{\BookTerm{T0283}{\csname BookOriginal@trTwistedComplex\endcsname}}
\expandafter\def\csname BookIndex@T0284\endcsname{niuqu@扭曲（twisted）}
\expandafter\let\csname BookOriginal@trTwisted\endcsname\trTwisted
\renewcommand{\trTwisted}{\BookTerm{T0284}{\csname BookOriginal@trTwisted\endcsname}}
\expandafter\def\csname BookIndex@T0285\endcsname{yingshezhui@映射锥（mapping cone）}
\expandafter\let\csname BookOriginal@trMappingCone\endcsname\trMappingCone
\renewcommand{\trMappingCone}{\BookTerm{T0285}{\csname BookOriginal@trMappingCone\endcsname}}
\expandafter\def\csname BookIndex@T0286\endcsname{xianzhishanglian@限制上链（bounding cochain）}
\expandafter\let\csname BookOriginal@trBoundingCochain\endcsname\trBoundingCochain
\renewcommand{\trBoundingCochain}{\BookTerm{T0286}{\csname BookOriginal@trBoundingCochain\endcsname}}
\expandafter\def\csname BookIndex@T0287\endcsname{zhui@锥（cone）}
\expandafter\let\csname BookOriginal@trCone\endcsname\trCone
\renewcommand{\trCone}{\BookTerm{T0287}{\csname BookOriginal@trCone\endcsname}}
\expandafter\def\csname BookIndex@T0288\endcsname{yiweihanzi@移位函子（shift functor）}
\expandafter\let\csname BookOriginal@trShiftFunctor\endcsname\trShiftFunctor
\renewcommand{\trShiftFunctor}{\BookTerm{T0288}{\csname BookOriginal@trShiftFunctor\endcsname}}
\expandafter\def\csname BookIndex@T0289\endcsname{yiwei@移位（suspension）}
\expandafter\let\csname BookOriginal@trSuspension\endcsname\trSuspension
\renewcommand{\trSuspension}{\BookTerm{T0289}{\csname BookOriginal@trSuspension\endcsname}}
\expandafter\def\csname BookIndex@T0290\endcsname{mitianqianru@米田嵌入（Yoneda embedding）}
\expandafter\let\csname BookOriginal@trYonedaEmbedding\endcsname\trYonedaEmbedding
\renewcommand{\trYonedaEmbedding}{\BookTerm{T0290}{\csname BookOriginal@trYonedaEmbedding\endcsname}}
\expandafter\def\csname BookIndex@T0291\endcsname{zhongshi@忠实（faithful）}
\expandafter\let\csname BookOriginal@trFaithful\endcsname\trFaithful
\renewcommand{\trFaithful}{\BookTerm{T0291}{\csname BookOriginal@trFaithful\endcsname}}
\expandafter\def\csname BookIndex@T0292\endcsname{man@满（full）}
\expandafter\let\csname BookOriginal@trFull\endcsname\trFull
\renewcommand{\trFull}{\BookTerm{T0292}{\csname BookOriginal@trFull\endcsname}}
\expandafter\def\csname BookIndex@T0293\endcsname{manqiezhongshi@满且忠实（fully faithful）}
\expandafter\let\csname BookOriginal@trFullyFaithful\endcsname\trFullyFaithful
\renewcommand{\trFullyFaithful}{\BookTerm{T0293}{\csname BookOriginal@trFullyFaithful\endcsname}}
\expandafter\def\csname BookIndex@T0294\endcsname{dengjia@等价（equivalence）}
\expandafter\let\csname BookOriginal@trEquivalence\endcsname\trEquivalence
\renewcommand{\trEquivalence}{\BookTerm{T0294}{\csname BookOriginal@trEquivalence\endcsname}}
\expandafter\def\csname BookIndex@T0295\endcsname{ziranbianhuan@自然变换（natural transformation）}
\expandafter\let\csname BookOriginal@trNaturalTransformation\endcsname\trNaturalTransformation
\renewcommand{\trNaturalTransformation}{\BookTerm{T0295}{\csname BookOriginal@trNaturalTransformation\endcsname}}
\expandafter\def\csname BookIndex@T0296\endcsname{daochufanchou@导出范畴（derived category）}
\expandafter\let\csname BookOriginal@trDerivedCategory\endcsname\trDerivedCategory
\renewcommand{\trDerivedCategory}{\BookTerm{T0296}{\csname BookOriginal@trDerivedCategory\endcsname}}
\expandafter\def\csname BookIndex@T0297\endcsname{huohexierdeshangtongdiao@霍赫希尔德上同调（Hochschild cohomology）}
\expandafter\let\csname BookOriginal@trHochschildCohomology\endcsname\trHochschildCohomology
\renewcommand{\trHochschildCohomology}{\BookTerm{T0297}{\csname BookOriginal@trHochschildCohomology\endcsname}}
\expandafter\def\csname BookIndex@T0298\endcsname{meixichengji@梅西乘积（Massey product）}
\expandafter\let\csname BookOriginal@trMasseyProduct\endcsname\trMasseyProduct
\renewcommand{\trMasseyProduct}{\BookTerm{T0298}{\csname BookOriginal@trMasseyProduct\endcsname}}
\expandafter\def\csname BookIndex@T0299\endcsname{xingshiweifentongpei@形式微分同胚（formal diffeomorphism）}
\expandafter\let\csname BookOriginal@trFormalDiffeomorphism\endcsname\trFormalDiffeomorphism
\renewcommand{\trFormalDiffeomorphism}{\BookTerm{T0299}{\csname BookOriginal@trFormalDiffeomorphism\endcsname}}
\expandafter\def\csname BookIndex@T0300\endcsname{zhuanyi@转移（transfer）}
\expandafter\let\csname BookOriginal@trTransfer\endcsname\trTransfer
\renewcommand{\trTransfer}{\BookTerm{T0300}{\csname BookOriginal@trTransfer\endcsname}}
\expandafter\def\csname BookIndex@T0301\endcsname{zhangai@障碍（obstruction）}
\expandafter\let\csname BookOriginal@trObstruction\endcsname\trObstruction
\renewcommand{\trObstruction}{\BookTerm{T0301}{\csname BookOriginal@trObstruction\endcsname}}
\expandafter\def\csname BookIndex@T0302\endcsname{puxulie@谱序列（spectral sequence）}
\expandafter\let\csname BookOriginal@trSpectralSequence\endcsname\trSpectralSequence
\renewcommand{\trSpectralSequence}{\BookTerm{T0302}{\csname BookOriginal@trSpectralSequence\endcsname}}
\expandafter\def\csname BookIndex@T0303\endcsname{zhenghedui@正合对（exact couple）}
\expandafter\let\csname BookOriginal@trExactCouple\endcsname\trExactCouple
\renewcommand{\trExactCouple}{\BookTerm{T0303}{\csname BookOriginal@trExactCouple\endcsname}}
\expandafter\def\csname BookIndex@T0304\endcsname{youxiang@有向（directed）}
\expandafter\let\csname BookOriginal@trDirected\endcsname\trDirected
\renewcommand{\trDirected}{\BookTerm{T0304}{\csname BookOriginal@trDirected\endcsname}}
\expandafter\def\csname BookIndex@T0305\endcsname{fanbian@反变（contravariant）}
\expandafter\let\csname BookOriginal@trContravariant\endcsname\trContravariant
\renewcommand{\trContravariant}{\BookTerm{T0305}{\csname BookOriginal@trContravariant\endcsname}}
\expandafter\def\csname BookIndex@T0306\endcsname{fenlie@分裂（splitting）}
\expandafter\let\csname BookOriginal@trSplitting\endcsname\trSplitting
\renewcommand{\trSplitting}{\BookTerm{T0306}{\csname BookOriginal@trSplitting\endcsname}}
\expandafter\def\csname BookIndex@T0307\endcsname{fenlieshengcheng@分裂生成（split generation）}
\expandafter\let\csname BookOriginal@trSplitGeneration\endcsname\trSplitGeneration
\renewcommand{\trSplitGeneration}{\BookTerm{T0307}{\csname BookOriginal@trSplitGeneration\endcsname}}
\expandafter\def\csname BookIndex@T0308\endcsname{fenliebi@分裂闭（split closed）}
\expandafter\let\csname BookOriginal@trSplitClosed\endcsname\trSplitClosed
\renewcommand{\trSplitClosed}{\BookTerm{T0308}{\csname BookOriginal@trSplitClosed\endcsname}}
\expandafter\def\csname BookIndex@T0309\endcsname{fenliebibao@分裂闭包（split closure）}
\expandafter\let\csname BookOriginal@trSplitClosure\endcsname\trSplitClosure
\renewcommand{\trSplitClosure}{\BookTerm{T0309}{\csname BookOriginal@trSplitClosure\endcsname}}
\expandafter\def\csname BookIndex@T0310\endcsname{bixian@闭弦（closed string）}
\expandafter\let\csname BookOriginal@trClosedString\endcsname\trClosedString
\renewcommand{\trClosedString}{\BookTerm{T0310}{\csname BookOriginal@trClosedString\endcsname}}
\expandafter\def\csname BookIndex@T0311\endcsname{kaixian@开弦（open string）}
\expandafter\let\csname BookOriginal@trOpenString\endcsname\trOpenString
\renewcommand{\trOpenString}{\BookTerm{T0311}{\csname BookOriginal@trOpenString\endcsname}}
\expandafter\def\csname BookIndex@T0312\endcsname{jiaxing@加性（additive）}
\expandafter\let\csname BookOriginal@trAdditive\endcsname\trAdditive
\renewcommand{\trAdditive}{\BookTerm{T0312}{\csname BookOriginal@trAdditive\endcsname}}
\expandafter\def\csname BookIndex@T0313\endcsname{bandan@半单（semisimple）}
\expandafter\let\csname BookOriginal@trSemisimple\endcsname\trSemisimple
\renewcommand{\trSemisimple}{\BookTerm{T0313}{\csname BookOriginal@trSemisimple\endcsname}}
\expandafter\def\csname BookIndex@T0314\endcsname{lingdiao@零调（acyclic）}
\expandafter\let\csname BookOriginal@trAcyclic\endcsname\trAcyclic
\renewcommand{\trAcyclic}{\BookTerm{T0314}{\csname BookOriginal@trAcyclic\endcsname}}
\expandafter\def\csname BookIndex@T0315\endcsname{shousuo@收缩（retract）}
\expandafter\let\csname BookOriginal@trRetract\endcsname\trRetract
\renewcommand{\trRetract}{\BookTerm{T0315}{\csname BookOriginal@trRetract\endcsname}}
\expandafter\def\csname BookIndex@T0316\endcsname{duoxianxing@多线性（multilinear）}
\expandafter\let\csname BookOriginal@trMultilinear\endcsname\trMultilinear
\renewcommand{\trMultilinear}{\BookTerm{T0316}{\csname BookOriginal@trMultilinear\endcsname}}
\expandafter\def\csname BookIndex@T0317\endcsname{fuhe@复合（composition）}
\expandafter\let\csname BookOriginal@trComposition\endcsname\trComposition
\renewcommand{\trComposition}{\BookTerm{T0317}{\csname BookOriginal@trComposition\endcsname}}
\expandafter\def\csname BookIndex@T0318\endcsname{ji@基（basis）}
\expandafter\let\csname BookOriginal@trBasis\endcsname\trBasis
\renewcommand{\trBasis}{\BookTerm{T0318}{\csname BookOriginal@trBasis\endcsname}}
\expandafter\def\csname BookIndex@T0319\endcsname{xishu@系数（coefficients）}
\expandafter\let\csname BookOriginal@trCoefficients\endcsname\trCoefficients
\renewcommand{\trCoefficients}{\BookTerm{T0319}{\csname BookOriginal@trCoefficients\endcsname}}
\expandafter\def\csname BookIndex@T0320\endcsname{xianxing@线性（linear）}
\expandafter\let\csname BookOriginal@trLinear\endcsname\trLinear
\renewcommand{\trLinear}{\BookTerm{T0320}{\csname BookOriginal@trLinear\endcsname}}
\expandafter\def\csname BookIndex@T0321\endcsname{shengufanchou@深谷范畴（Fukaya category）}
\expandafter\let\csname BookOriginal@trFukayaCategory\endcsname\trFukayaCategory
\renewcommand{\trFukayaCategory}{\BookTerm{T0321}{\csname BookOriginal@trFukayaCategory\endcsname}}
\expandafter\def\csname BookIndex@T0322\endcsname{mo@膜（brane）}
\expandafter\let\csname BookOriginal@trBrane\endcsname\trBrane
\renewcommand{\trBrane}{\BookTerm{T0322}{\csname BookOriginal@trBrane\endcsname}}
\expandafter\def\csname BookIndex@T0323\endcsname{jingquelagelangrimo@精确拉格朗日膜（exact Lagrangian brane）}
\expandafter\let\csname BookOriginal@trExactLagrangianBrane\endcsname\trExactLagrangianBrane
\renewcommand{\trExactLagrangianBrane}{\BookTerm{T0323}{\csname BookOriginal@trExactLagrangianBrane\endcsname}}
\expandafter\def\csname BookIndex@T0324\endcsname{fenci@分次（grading）}
\expandafter\let\csname BookOriginal@trGrading\endcsname\trGrading
\renewcommand{\trGrading}{\BookTerm{T0324}{\csname BookOriginal@trGrading\endcsname}}
\expandafter\def\csname BookIndex@T0325\endcsname{pinjiegou@Pin~结构（pin structure）}
\expandafter\let\csname BookOriginal@trPinStructure\endcsname\trPinStructure
\renewcommand{\trPinStructure}{\BookTerm{T0325}{\csname BookOriginal@trPinStructure\endcsname}}
\expandafter\def\csname BookIndex@T0326\endcsname{xinglieshixian@行列式线（determinant line）}
\expandafter\let\csname BookOriginal@trDeterminantLine\endcsname\trDeterminantLine
\renewcommand{\trDeterminantLine}{\BookTerm{T0326}{\csname BookOriginal@trDeterminantLine\endcsname}}
\expandafter\def\csname BookIndex@T0327\endcsname{duobianxing@多边形（polygon）}
\expandafter\let\csname BookOriginal@trPolygon\endcsname\trPolygon
\renewcommand{\trPolygon}{\BookTerm{T0327}{\csname BookOriginal@trPolygon\endcsname}}
\expandafter\def\csname BookIndex@T0328\endcsname{bianjietiaojian@边界条件（boundary condition）}
\expandafter\let\csname BookOriginal@trBoundaryCondition\endcsname\trBoundaryCondition
\renewcommand{\trBoundaryCondition}{\BookTerm{T0328}{\csname BookOriginal@trBoundaryCondition\endcsname}}
\expandafter\def\csname BookIndex@T0329\endcsname{wendingyingshe@稳定映射（stable map）}
\expandafter\let\csname BookOriginal@trStableMap\endcsname\trStableMap
\renewcommand{\trStableMap}{\BookTerm{T0329}{\csname BookOriginal@trStableMap\endcsname}}
\expandafter\def\csname BookIndex@T0330\endcsname{kuige@亏格（genus）}
\expandafter\let\csname BookOriginal@trGenus\endcsname\trGenus
\renewcommand{\trGenus}{\BookTerm{T0330}{\csname BookOriginal@trGenus\endcsname}}
\expandafter\def\csname BookIndex@T0331\endcsname{biaojidian@标记点（marked point）}
\expandafter\let\csname BookOriginal@trMarkedPoint\endcsname\trMarkedPoint
\renewcommand{\trMarkedPoint}{\BookTerm{T0331}{\csname BookOriginal@trMarkedPoint\endcsname}}
\expandafter\def\csname BookIndex@T0332\endcsname{daizhuangduan@带状端（striplike end）}
\expandafter\let\csname BookOriginal@trStriplikeEnd\endcsname\trStriplikeEnd
\renewcommand{\trStriplikeEnd}{\BookTerm{T0332}{\csname BookOriginal@trStriplikeEnd\endcsname}}
\expandafter\def\csname BookIndex@T0333\endcsname{fuzhiyingshe@赋值映射（evaluation map）}
\expandafter\let\csname BookOriginal@trEvaluationMap\endcsname\trEvaluationMap
\renewcommand{\trEvaluationMap}{\BookTerm{T0333}{\csname BookOriginal@trEvaluationMap\endcsname}}
\expandafter\def\csname BookIndex@T0334\endcsname{xianweiji@纤维积（fibre product）}
\expandafter\let\csname BookOriginal@trFibreProduct\endcsname\trFibreProduct
\renewcommand{\trFibreProduct}{\BookTerm{T0334}{\csname BookOriginal@trFibreProduct\endcsname}}
\expandafter\def\csname BookIndex@T0335\endcsname{lahui@拉回（pullback）}
\expandafter\let\csname BookOriginal@trPullback\endcsname\trPullback
\renewcommand{\trPullback}{\BookTerm{T0335}{\csname BookOriginal@trPullback\endcsname}}
\expandafter\def\csname BookIndex@T0336\endcsname{tuiqian@推前（pushforward）}
\expandafter\let\csname BookOriginal@trPushforward\endcsname\trPushforward
\renewcommand{\trPushforward}{\BookTerm{T0336}{\csname BookOriginal@trPushforward\endcsname}}
\expandafter\def\csname BookIndex@T0337\endcsname{pingtan@平坦（flat）}
\expandafter\let\csname BookOriginal@trFlat\endcsname\trFlat
\renewcommand{\trFlat}{\BookTerm{T0337}{\csname BookOriginal@trFlat\endcsname}}
\expandafter\def\csname BookIndex@T0338\endcsname{lianluo@联络（connection）}
\expandafter\let\csname BookOriginal@trConnection\endcsname\trConnection
\renewcommand{\trConnection}{\BookTerm{T0338}{\csname BookOriginal@trConnection\endcsname}}
\expandafter\def\csname BookIndex@T0339\endcsname{dengbian@等变（equivariant）}
\expandafter\let\csname BookOriginal@trEquivariant\endcsname\trEquivariant
\renewcommand{\trEquivariant}{\BookTerm{T0339}{\csname BookOriginal@trEquivariant\endcsname}}
\expandafter\def\csname BookIndex@T0340\endcsname{duihe@对合（involution）}
\expandafter\let\csname BookOriginal@trInvolution\endcsname\trInvolution
\renewcommand{\trInvolution}{\BookTerm{T0340}{\csname BookOriginal@trInvolution\endcsname}}
\expandafter\def\csname BookIndex@T0341\endcsname{shang@商（quotient）}
\expandafter\let\csname BookOriginal@trQuotient\endcsname\trQuotient
\renewcommand{\trQuotient}{\BookTerm{T0341}{\csname BookOriginal@trQuotient\endcsname}}
\expandafter\def\csname BookIndex@T0342\endcsname{juerejixian@绝热极限（adiabatic limit）}
\expandafter\let\csname BookOriginal@trAdiabaticLimit\endcsname\trAdiabaticLimit
\renewcommand{\trAdiabaticLimit}{\BookTerm{T0342}{\csname BookOriginal@trAdiabaticLimit\endcsname}}
\expandafter\def\csname BookIndex@T0343\endcsname{zhongbiaodu@重标度（rescaling）}
\expandafter\let\csname BookOriginal@trRescaling\endcsname\trRescaling
\renewcommand{\trRescaling}{\BookTerm{T0343}{\csname BookOriginal@trRescaling\endcsname}}
\expandafter\def\csname BookIndex@T0344\endcsname{peidui@配对（pairing）}
\expandafter\let\csname BookOriginal@trPairing\endcsname\trPairing
\renewcommand{\trPairing}{\BookTerm{T0344}{\csname BookOriginal@trPairing\endcsname}}
\expandafter\def\csname BookIndex@T0345\endcsname{tianchong@填充（filling）}
\expandafter\let\csname BookOriginal@trFilling\endcsname\trFilling
\renewcommand{\trFilling}{\BookTerm{T0345}{\csname BookOriginal@trFilling\endcsname}}
\expandafter\def\csname BookIndex@T0346\endcsname{quanchunyuanpan@全纯圆盘（holomorphic disc）}
\expandafter\let\csname BookOriginal@trHolomorphicDisc\endcsname\trHolomorphicDisc
\renewcommand{\trHolomorphicDisc}{\BookTerm{T0346}{\csname BookOriginal@trHolomorphicDisc\endcsname}}
\expandafter\def\csname BookIndex@T0347\endcsname{quanchunduobianxing@全纯多边形（holomorphic polygon）}
\expandafter\let\csname BookOriginal@trHolomorphicPolygon\endcsname\trHolomorphicPolygon
\renewcommand{\trHolomorphicPolygon}{\BookTerm{T0347}{\csname BookOriginal@trHolomorphicPolygon\endcsname}}
\expandafter\def\csname BookIndex@T0348\endcsname{moersibote@莫尔斯--博特（Morse Bott）}
\expandafter\let\csname BookOriginal@trMorseBott\endcsname\trMorseBott
\renewcommand{\trMorseBott}{\BookTerm{T0348}{\csname BookOriginal@trMorseBott\endcsname}}
\expandafter\def\csname BookIndex@T0349\endcsname{fuhelu@复合律（composition law）}
\expandafter\let\csname BookOriginal@trCompositionLaw\endcsname\trCompositionLaw
\renewcommand{\trCompositionLaw}{\BookTerm{T0349}{\csname BookOriginal@trCompositionLaw\endcsname}}
\expandafter\def\csname BookIndex@T0350\endcsname{sanjiaochengji@三角乘积（triangle product）}
\expandafter\let\csname BookOriginal@trTriangleProduct\endcsname\trTriangleProduct
\renewcommand{\trTriangleProduct}{\BookTerm{T0350}{\csname BookOriginal@trTriangleProduct\endcsname}}
\expandafter\def\csname BookIndex@T0351\endcsname{feiqici@非齐次（inhomogeneous）}
\expandafter\let\csname BookOriginal@trInhomogeneous\endcsname\trInhomogeneous
\renewcommand{\trInhomogeneous}{\BookTerm{T0351}{\csname BookOriginal@trInhomogeneous\endcsname}}
\expandafter\def\csname BookIndex@T0352\endcsname{weiquanchunyingshe@伪全纯映射（pseudoholomorphic map）}
\expandafter\let\csname BookOriginal@trPseudoholomorphicMap\endcsname\trPseudoholomorphicMap
\renewcommand{\trPseudoholomorphicMap}{\BookTerm{T0352}{\csname BookOriginal@trPseudoholomorphicMap\endcsname}}
\expandafter\def\csname BookIndex@T0353\endcsname{fuluoershuju@弗洛尔数据（Floer datum）}
\expandafter\let\csname BookOriginal@trFloerDatum\endcsname\trFloerDatum
\renewcommand{\trFloerDatum}{\BookTerm{T0353}{\csname BookOriginal@trFloerDatum\endcsname}}
\expandafter\def\csname BookIndex@T0354\endcsname{raodongshuju@扰动数据（perturbation datum）}
\expandafter\let\csname BookOriginal@trPerturbationDatum\endcsname\trPerturbationDatum
\renewcommand{\trPerturbationDatum}{\BookTerm{T0354}{\csname BookOriginal@trPerturbationDatum\endcsname}}
\expandafter\def\csname BookIndex@T0355\endcsname{jianjin@渐近（asymptotic）}
\expandafter\let\csname BookOriginal@trAsymptotic\endcsname\trAsymptotic
\renewcommand{\trAsymptotic}{\BookTerm{T0355}{\csname BookOriginal@trAsymptotic\endcsname}}
\expandafter\def\csname BookIndex@T0356\endcsname{lujingkongjian@路径空间（path space）}
\expandafter\let\csname BookOriginal@trPathSpace\endcsname\trPathSpace
\renewcommand{\trPathSpace}{\BookTerm{T0356}{\csname BookOriginal@trPathSpace\endcsname}}
\expandafter\def\csname BookIndex@T0357\endcsname{qulu@曲率（curvature）}
\expandafter\let\csname BookOriginal@trCurvature\endcsname\trCurvature
\renewcommand{\trCurvature}{\BookTerm{T0357}{\csname BookOriginal@trCurvature\endcsname}}
\expandafter\def\csname BookIndex@T0358\endcsname{gongxing@共形（conformal）}
\expandafter\let\csname BookOriginal@trConformal\endcsname\trConformal
\renewcommand{\trConformal}{\BookTerm{T0358}{\csname BookOriginal@trConformal\endcsname}}
\expandafter\def\csname BookIndex@T0359\endcsname{pikalaifuxiecililun@皮卡--莱夫谢茨理论（Picard–Lefschetz theory）}
\expandafter\let\csname BookOriginal@trPicardLefschetzTheory\endcsname\trPicardLefschetzTheory
\renewcommand{\trPicardLefschetzTheory}{\BookTerm{T0359}{\csname BookOriginal@trPicardLefschetzTheory\endcsname}}
\expandafter\def\csname BookIndex@T0360\endcsname{laifuxiecixianweihua@莱夫谢茨纤维化（Lefschetz fibration）}
\expandafter\let\csname BookOriginal@trLefschetzFibration\endcsname\trLefschetzFibration
\renewcommand{\trLefschetzFibration}{\BookTerm{T0360}{\csname BookOriginal@trLefschetzFibration\endcsname}}
\expandafter\def\csname BookIndex@T0361\endcsname{laifuxieciqianbi@莱夫谢茨铅笔（Lefschetz pencil）}
\expandafter\let\csname BookOriginal@trLefschetzPencil\endcsname\trLefschetzPencil
\renewcommand{\trLefschetzPencil}{\BookTerm{T0361}{\csname BookOriginal@trLefschetzPencil\endcsname}}
\expandafter\def\csname BookIndex@T0362\endcsname{laifuxiecizhen@莱夫谢茨针（Lefschetz thimble）}
\expandafter\let\csname BookOriginal@trLefschetzThimble\endcsname\trLefschetzThimble
\renewcommand{\trLefschetzThimble}{\BookTerm{T0362}{\csname BookOriginal@trLefschetzThimble\endcsname}}
\expandafter\def\csname BookIndex@T0363\endcsname{xiaomeibilian@消没闭链（vanishing cycle）}
\expandafter\let\csname BookOriginal@trVanishingCycle\endcsname\trVanishingCycle
\renewcommand{\trVanishingCycle}{\BookTerm{T0363}{\csname BookOriginal@trVanishingCycle\endcsname}}
\expandafter\def\csname BookIndex@T0364\endcsname{xiaomeilujing@消没路径（vanishing path）}
\expandafter\let\csname BookOriginal@trVanishingPath\endcsname\trVanishingPath
\renewcommand{\trVanishingPath}{\BookTerm{T0364}{\csname BookOriginal@trVanishingPath\endcsname}}
\expandafter\def\csname BookIndex@T0365\endcsname{pipeibilian@匹配闭链（matching cycle）}
\expandafter\let\csname BookOriginal@trMatchingCycle\endcsname\trMatchingCycle
\renewcommand{\trMatchingCycle}{\BookTerm{T0365}{\csname BookOriginal@trMatchingCycle\endcsname}}
\expandafter\def\csname BookIndex@T0366\endcsname{pipeilujing@匹配路径（matching path）}
\expandafter\let\csname BookOriginal@trMatchingPath\endcsname\trMatchingPath
\renewcommand{\trMatchingPath}{\BookTerm{T0366}{\csname BookOriginal@trMatchingPath\endcsname}}
\expandafter\def\csname BookIndex@T0367\endcsname{miernuoxianwei@米尔诺纤维（Milnor fibre）}
\expandafter\let\csname BookOriginal@trMilnorFibre\endcsname\trMilnorFibre
\renewcommand{\trMilnorFibre}{\BookTerm{T0367}{\csname BookOriginal@trMilnorFibre\endcsname}}
\expandafter\def\csname BookIndex@T0368\endcsname{quankongjian@全空间（total space）}
\expandafter\let\csname BookOriginal@trTotalSpace\endcsname\trTotalSpace
\renewcommand{\trTotalSpace}{\BookTerm{T0368}{\csname BookOriginal@trTotalSpace\endcsname}}
\expandafter\def\csname BookIndex@T0369\endcsname{qiyixianwei@奇异纤维（singular fibre）}
\expandafter\let\csname BookOriginal@trSingularFibre\endcsname\trSingularFibre
\renewcommand{\trSingularFibre}{\BookTerm{T0369}{\csname BookOriginal@trSingularFibre\endcsname}}
\expandafter\def\csname BookIndex@T0370\endcsname{deenniuzhuan@德恩扭转（Dehn twist）}
\expandafter\let\csname BookOriginal@trDehnTwist\endcsname\trDehnTwist
\renewcommand{\trDehnTwist}{\BookTerm{T0370}{\csname BookOriginal@trDehnTwist\endcsname}}
\expandafter\def\csname BookIndex@T0371\endcsname{niuzhuanhanzi@扭转函子（twist functor）}
\expandafter\let\csname BookOriginal@trTwistFunctor\endcsname\trTwistFunctor
\renewcommand{\trTwistFunctor}{\BookTerm{T0371}{\csname BookOriginal@trTwistFunctor\endcsname}}
\expandafter\def\csname BookIndex@T0372\endcsname{bian@辫（braid）}
\expandafter\let\csname BookOriginal@trBraid\endcsname\trBraid
\renewcommand{\trBraid}{\BookTerm{T0372}{\csname BookOriginal@trBraid\endcsname}}
\expandafter\def\csname BookIndex@T0373\endcsname{bianqun@辫群（braid group）}
\expandafter\let\csname BookOriginal@trBraidGroup\endcsname\trBraidGroup
\renewcommand{\trBraidGroup}{\BookTerm{T0373}{\csname BookOriginal@trBraidGroup\endcsname}}
\expandafter\def\csname BookIndex@T0374\endcsname{daishucu@代数簇（algebraic variety）}
\expandafter\let\csname BookOriginal@trAlgebraicVariety\endcsname\trAlgebraicVariety
\renewcommand{\trAlgebraicVariety}{\BookTerm{T0374}{\csname BookOriginal@trAlgebraicVariety\endcsname}}
\expandafter\def\csname BookIndex@T0375\endcsname{sheyingcu@射影簇（projective variety）}
\expandafter\let\csname BookOriginal@trProjectiveVariety\endcsname\trProjectiveVariety
\renewcommand{\trProjectiveVariety}{\BookTerm{T0375}{\csname BookOriginal@trProjectiveVariety\endcsname}}
\expandafter\def\csname BookIndex@T0376\endcsname{fangshecu@仿射簇（affine variety）}
\expandafter\let\csname BookOriginal@trAffineVariety\endcsname\trAffineVariety
\renewcommand{\trAffineVariety}{\BookTerm{T0376}{\csname BookOriginal@trAffineVariety\endcsname}}
\expandafter\def\csname BookIndex@T0377\endcsname{chaopingmianjiemian@超平面截面（hyperplane section）}
\expandafter\let\csname BookOriginal@trHyperplaneSection\endcsname\trHyperplaneSection
\renewcommand{\trHyperplaneSection}{\BookTerm{T0377}{\csname BookOriginal@trHyperplaneSection\endcsname}}
\expandafter\def\csname BookIndex@T0378\endcsname{fengfu@丰富（ample）}
\expandafter\let\csname BookOriginal@trAmple\endcsname\trAmple
\renewcommand{\trAmple}{\BookTerm{T0378}{\csname BookOriginal@trAmple\endcsname}}
\expandafter\def\csname BookIndex@T0379\endcsname{fenqifugai@分歧覆盖（branched cover）}
\expandafter\let\csname BookOriginal@trBranchedCover\endcsname\trBranchedCover
\renewcommand{\trBranchedCover}{\BookTerm{T0379}{\csname BookOriginal@trBranchedCover\endcsname}}
\expandafter\def\csname BookIndex@T0380\endcsname{fenqidian@分歧点（branch point）}
\expandafter\let\csname BookOriginal@trBranchPoint\endcsname\trBranchPoint
\renewcommand{\trBranchPoint}{\BookTerm{T0380}{\csname BookOriginal@trBranchPoint\endcsname}}
\expandafter\def\csname BookIndex@T0381\endcsname{shuiping@水平（horizontal）}
\expandafter\let\csname BookOriginal@trHorizontal\endcsname\trHorizontal
\renewcommand{\trHorizontal}{\BookTerm{T0381}{\csname BookOriginal@trHorizontal\endcsname}}
\expandafter\def\csname BookIndex@T0382\endcsname{shuzhi@竖直（vertical）}
\expandafter\let\csname BookOriginal@trVertical\endcsname\trVertical
\renewcommand{\trVertical}{\BookTerm{T0382}{\csname BookOriginal@trVertical\endcsname}}
\expandafter\def\csname BookIndex@T0383\endcsname{xianweihua@纤维化（fibration）}
\expandafter\let\csname BookOriginal@trFibration\endcsname\trFibration
\renewcommand{\trFibration}{\BookTerm{T0383}{\csname BookOriginal@trFibration\endcsname}}
\expandafter\def\csname BookIndex@T0384\endcsname{yuqiexianwei@余切纤维（cotangent fibre）}
\expandafter\let\csname BookOriginal@trCotangentFibre\endcsname\trCotangentFibre
\renewcommand{\trCotangentFibre}{\BookTerm{T0384}{\csname BookOriginal@trCotangentFibre\endcsname}}
\expandafter\def\csname BookIndex@T0385\endcsname{wanbeihua@完备化（completion）}
\expandafter\let\csname BookOriginal@trCompletion\endcsname\trCompletion
\renewcommand{\trCompletion}{\BookTerm{T0385}{\csname BookOriginal@trCompletion\endcsname}}
\expandafter\def\csname BookIndex@T0386\endcsname{zu@族（family）}
\expandafter\let\csname BookOriginal@trFamily\endcsname\trFamily
\renewcommand{\trFamily}{\BookTerm{T0386}{\csname BookOriginal@trFamily\endcsname}}
\expandafter\def\csname BookIndex@T0387\endcsname{zhuxianwei@逐纤维（fibrewise）}
\expandafter\let\csname BookOriginal@trFibrewise\endcsname\trFibrewise
\renewcommand{\trFibrewise}{\BookTerm{T0387}{\csname BookOriginal@trFibrewise\endcsname}}
\expandafter\def\csname BookIndex@T0388\endcsname{jiegouqun@结构群（structure group）}
\expandafter\let\csname BookOriginal@trStructureGroup\endcsname\trStructureGroup
\renewcommand{\trStructureGroup}{\BookTerm{T0388}{\csname BookOriginal@trStructureGroup\endcsname}}
\expandafter\def\csname BookIndex@T0389\endcsname{quanshi@全实（totally real）}
\expandafter\let\csname BookOriginal@trTotallyReal\endcsname\trTotallyReal
\renewcommand{\trTotallyReal}{\BookTerm{T0389}{\csname BookOriginal@trTotallyReal\endcsname}}
\expandafter\def\csname BookIndex@T0390\endcsname{jiaquan@加权（weighted）}
\expandafter\let\csname BookOriginal@trWeighted\endcsname\trWeighted
\renewcommand{\trWeighted}{\BookTerm{T0390}{\csname BookOriginal@trWeighted\endcsname}}
\expandafter\def\csname BookIndex@T0391\endcsname{quanzhong@权重（weight）}
\expandafter\let\csname BookOriginal@trWeight\endcsname\trWeight
\renewcommand{\trWeight}{\BookTerm{T0391}{\csname BookOriginal@trWeight\endcsname}}
\expandafter\def\csname BookIndex@T0392\endcsname{dingyiyu@定义域（domain）}
\expandafter\let\csname BookOriginal@trDomain\endcsname\trDomain
\renewcommand{\trDomain}{\BookTerm{T0392}{\csname BookOriginal@trDomain\endcsname}}
\expandafter\def\csname BookIndex@T0393\endcsname{jixiaofan@极小泛（miniversal）}
\expandafter\let\csname BookOriginal@trMiniversal\endcsname\trMiniversal
\renewcommand{\trMiniversal}{\BookTerm{T0393}{\csname BookOriginal@trMiniversal\endcsname}}
\expandafter\def\csname BookIndex@T0394\endcsname{wanyou@万有（universal）}
\expandafter\let\csname BookOriginal@trUniversal\endcsname\trUniversal
\renewcommand{\trUniversal}{\BookTerm{T0394}{\csname BookOriginal@trUniversal\endcsname}}
\expandafter\def\csname BookIndex@T0395\endcsname{daibiaojiyuanpan@带标记圆盘（pointed disc）}
\expandafter\let\csname BookOriginal@trPointedDisc\endcsname\trPointedDisc
\renewcommand{\trPointedDisc}{\BookTerm{T0395}{\csname BookOriginal@trPointedDisc\endcsname}}
\expandafter\def\csname BookIndex@T0396\endcsname{xunhuancixu@循环次序（cyclic order）}
\expandafter\let\csname BookOriginal@trCyclicOrder\endcsname\trCyclicOrder
\renewcommand{\trCyclicOrder}{\BookTerm{T0396}{\csname BookOriginal@trCyclicOrder\endcsname}}
\expandafter\def\csname BookIndex@T0397\endcsname{neidian@内点（interior point）}
\expandafter\let\csname BookOriginal@trInteriorPoint\endcsname\trInteriorPoint
\renewcommand{\trInteriorPoint}{\BookTerm{T0397}{\csname BookOriginal@trInteriorPoint\endcsname}}
\expandafter\def\csname BookIndex@T0398\endcsname{yinlihoudai@引力后代（gravitational descendant）}
\expandafter\let\csname BookOriginal@trGravitationalDescendant\endcsname\trGravitationalDescendant
\renewcommand{\trGravitationalDescendant}{\BookTerm{T0398}{\csname BookOriginal@trGravitationalDescendant\endcsname}}
\expandafter\def\csname BookIndex@T0399\endcsname{gouxingkongjian@构型空间（configuration space）}
\expandafter\let\csname BookOriginal@trConfigurationSpace\endcsname\trConfigurationSpace
\renewcommand{\trConfigurationSpace}{\BookTerm{T0399}{\csname BookOriginal@trConfigurationSpace\endcsname}}
\expandafter\def\csname BookIndex@T0400\endcsname{danzhihua@单值化（uniformization）}
\expandafter\let\csname BookOriginal@trUniformization\endcsname\trUniformization
\renewcommand{\trUniformization}{\BookTerm{T0400}{\csname BookOriginal@trUniformization\endcsname}}
\expandafter\def\csname BookIndex@T0401\endcsname{fenleiyingshe@分类映射（classifying map）}
\expandafter\let\csname BookOriginal@trClassifyingMap\endcsname\trClassifyingMap
\renewcommand{\trClassifyingMap}{\BookTerm{T0401}{\csname BookOriginal@trClassifyingMap\endcsname}}
\expandafter\def\csname BookIndex@T0402\endcsname{shu@树（tree）}
\expandafter\let\csname BookOriginal@trTree\endcsname\trTree
\renewcommand{\trTree}{\BookTerm{T0402}{\csname BookOriginal@trTree\endcsname}}
\expandafter\def\csname BookIndex@T0403\endcsname{yeshu@叶树（leafed tree）}
\expandafter\let\csname BookOriginal@trLeafedTree\endcsname\trLeafedTree
\renewcommand{\trLeafedTree}{\BookTerm{T0403}{\csname BookOriginal@trLeafedTree\endcsname}}
\expandafter\def\csname BookIndex@T0404\endcsname{ye@叶（leaf）}
\expandafter\let\csname BookOriginal@trLeaf\endcsname\trLeaf
\renewcommand{\trLeaf}{\BookTerm{T0404}{\csname BookOriginal@trLeaf\endcsname}}
\expandafter\def\csname BookIndex@T0405\endcsname{gen@根（root）}
\expandafter\let\csname BookOriginal@trRoot\endcsname\trRoot
\renewcommand{\trRoot}{\BookTerm{T0405}{\csname BookOriginal@trRoot\endcsname}}
\expandafter\def\csname BookIndex@T0406\endcsname{dingdian@顶点（vertex）}
\expandafter\let\csname BookOriginal@trVertex\endcsname\trVertex
\renewcommand{\trVertex}{\BookTerm{T0406}{\csname BookOriginal@trVertex\endcsname}}
\expandafter\def\csname BookIndex@T0407\endcsname{bian@边（edge）}
\expandafter\let\csname BookOriginal@trEdge\endcsname\trEdge
\renewcommand{\trEdge}{\BookTerm{T0407}{\csname BookOriginal@trEdge\endcsname}}
\expandafter\def\csname BookIndex@T0408\endcsname{qi@旗（flag）}
\expandafter\let\csname BookOriginal@trFlag\endcsname\trFlag
\renewcommand{\trFlag}{\BookTerm{T0408}{\csname BookOriginal@trFlag\endcsname}}
\expandafter\def\csname BookIndex@T0409\endcsname{jia@价（valency）}
\expandafter\let\csname BookOriginal@trValency\endcsname\trValency
\renewcommand{\trValency}{\BookTerm{T0409}{\csname BookOriginal@trValency\endcsname}}
\expandafter\def\csname BookIndex@T0410\endcsname{waibubian@外部边（exterior edge）}
\expandafter\let\csname BookOriginal@trExteriorEdge\endcsname\trExteriorEdge
\renewcommand{\trExteriorEdge}{\BookTerm{T0410}{\csname BookOriginal@trExteriorEdge\endcsname}}
\expandafter\def\csname BookIndex@T0411\endcsname{neibubian@内部边（interior edge）}
\expandafter\let\csname BookOriginal@trInteriorEdge\endcsname\trInteriorEdge
\renewcommand{\trInteriorEdge}{\BookTerm{T0411}{\csname BookOriginal@trInteriorEdge\endcsname}}
\expandafter\def\csname BookIndex@T0412\endcsname{neibuqi@内部旗（interior flag）}
\expandafter\let\csname BookOriginal@trInteriorFlag\endcsname\trInteriorFlag
\renewcommand{\trInteriorFlag}{\BookTerm{T0412}{\csname BookOriginal@trInteriorFlag\endcsname}}
\expandafter\def\csname BookIndex@T0413\endcsname{nishizhen@逆时针（anticlockwise）}
\expandafter\let\csname BookOriginal@trAnticlockwise\endcsname\trAnticlockwise
\renewcommand{\trAnticlockwise}{\BookTerm{T0413}{\csname BookOriginal@trAnticlockwise\endcsname}}
\expandafter\def\csname BookIndex@T0414\endcsname{xinjiegou@辛结构（symplectic structure）}
\expandafter\let\csname BookOriginal@trSymplecticStructure\endcsname\trSymplecticStructure
\renewcommand{\trSymplecticStructure}{\BookTerm{T0414}{\csname BookOriginal@trSymplecticStructure\endcsname}}
\expandafter\def\csname BookIndex@T0415\endcsname{jixianshuju@极限数据（limiting datum）}
\expandafter\let\csname BookOriginal@trLimitingDatum\endcsname\trLimitingDatum
\renewcommand{\trLimitingDatum}{\BookTerm{T0415}{\csname BookOriginal@trLimitingDatum\endcsname}}
\expandafter\def\csname BookIndex@T0416\endcsname{kerongxu@可容许（admissible）}
\expandafter\let\csname BookOriginal@trAdmissible\endcsname\trAdmissible
\renewcommand{\trAdmissible}{\BookTerm{T0416}{\csname BookOriginal@trAdmissible\endcsname}}
\expandafter\def\csname BookIndex@T0417\endcsname{xiebianweifen@协变微分（covariant differentiation）}
\expandafter\let\csname BookOriginal@trCovariantDifferentiation\endcsname\trCovariantDifferentiation
\renewcommand{\trCovariantDifferentiation}{\BookTerm{T0417}{\csname BookOriginal@trCovariantDifferentiation\endcsname}}
\expandafter\def\csname BookIndex@T0418\endcsname{jinxing@紧性（compactness）}
\expandafter\let\csname BookOriginal@trCompactness\endcsname\trCompactness
\renewcommand{\trCompactness}{\BookTerm{T0418}{\csname BookOriginal@trCompactness\endcsname}}
\expandafter\def\csname BookIndex@T0419\endcsname{fanxing@泛性（versality）}
\expandafter\let\csname BookOriginal@trVersality\endcsname\trVersality
\renewcommand{\trVersality}{\BookTerm{T0419}{\csname BookOriginal@trVersality\endcsname}}
\expandafter\def\csname BookIndex@T0420\endcsname{yidongbianjietiaojian@移动边界条件（moving boundary condition）}
\expandafter\let\csname BookOriginal@trMovingBoundaryCondition\endcsname\trMovingBoundaryCondition
\renewcommand{\trMovingBoundaryCondition}{\BookTerm{T0420}{\csname BookOriginal@trMovingBoundaryCondition\endcsname}}
\expandafter\def\csname BookIndex@T0421\endcsname{shangji@上积（cup product）}
\expandafter\let\csname BookOriginal@trCupProduct\endcsname\trCupProduct
\renewcommand{\trCupProduct}{\BookTerm{T0421}{\csname BookOriginal@trCupProduct\endcsname}}
\expandafter\def\csname BookIndex@T0422\endcsname{bianjieliantonghe@边界连通和（boundary connected sum）}
\expandafter\let\csname BookOriginal@trBoundaryConnectedSum\endcsname\trBoundaryConnectedSum
\renewcommand{\trBoundaryConnectedSum}{\BookTerm{T0422}{\csname BookOriginal@trBoundaryConnectedSum\endcsname}}
\expandafter\def\csname BookIndex@T0423\endcsname{zhanheyingshe@粘合映射（gluing map）}
\expandafter\let\csname BookOriginal@trGluingMap\endcsname\trGluingMap
\renewcommand{\trGluingMap}{\BookTerm{T0423}{\csname BookOriginal@trGluingMap\endcsname}}
\expandafter\def\csname BookIndex@T0424\endcsname{zhanhecanshu@粘合参数（gluing parameter）}
\expandafter\let\csname BookOriginal@trGluingParameter\endcsname\trGluingParameter
\renewcommand{\trGluingParameter}{\BookTerm{T0424}{\csname BookOriginal@trGluingParameter\endcsname}}
\expandafter\def\csname BookIndex@T0425\endcsname{zhanhegouzao@粘合构造（gluing construction）}
\expandafter\let\csname BookOriginal@trGluingConstruction\endcsname\trGluingConstruction
\renewcommand{\trGluingConstruction}{\BookTerm{T0425}{\csname BookOriginal@trGluingConstruction\endcsname}}
\expandafter\def\csname BookIndex@T0426\endcsname{zhanhezhangdu@粘合长度（gluing length）}
\expandafter\let\csname BookOriginal@trGluingLength\endcsname\trGluingLength
\renewcommand{\trGluingLength}{\BookTerm{T0426}{\csname BookOriginal@trGluingLength\endcsname}}
\expandafter\def\csname BookIndex@T0427\endcsname{bufenjinhua@部分紧化（partial compactification）}
\expandafter\let\csname BookOriginal@trPartialCompactification\endcsname\trPartialCompactification
\renewcommand{\trPartialCompactification}{\BookTerm{T0427}{\csname BookOriginal@trPartialCompactification\endcsname}}
\expandafter\def\csname BookIndex@T0428\endcsname{guanghuajiegou@光滑结构（smooth structure）}
\expandafter\let\csname BookOriginal@trSmoothStructure\endcsname\trSmoothStructure
\renewcommand{\trSmoothStructure}{\BookTerm{T0428}{\csname BookOriginal@trSmoothStructure\endcsname}}
\expandafter\def\csname BookIndex@T0429\endcsname{yizhi@一致（consistent）}
\expandafter\let\csname BookOriginal@trConsistent\endcsname\trConsistent
\renewcommand{\trConsistent}{\BookTerm{T0429}{\csname BookOriginal@trConsistent\endcsname}}
\expandafter\def\csname BookIndex@T0430\endcsname{yizhixing@一致性（consistency）}
\expandafter\let\csname BookOriginal@trConsistency\endcsname\trConsistency
\renewcommand{\trConsistency}{\BookTerm{T0430}{\csname BookOriginal@trConsistency\endcsname}}
\expandafter\def\csname BookIndex@T0431\endcsname{jubuweifentongpei@局部微分同胚（local diffeomorphism）}
\expandafter\let\csname BookOriginal@trLocalDiffeomorphism\endcsname\trLocalDiffeomorphism
\renewcommand{\trLocalDiffeomorphism}{\BookTerm{T0431}{\csname BookOriginal@trLocalDiffeomorphism\endcsname}}
\expandafter\def\csname BookIndex@T0432\endcsname{kaiqianru@开嵌入（open embedding）}
\expandafter\let\csname BookOriginal@trOpenEmbedding\endcsname\trOpenEmbedding
\renewcommand{\trOpenEmbedding}{\BookTerm{T0432}{\csname BookOriginal@trOpenEmbedding\endcsname}}
\expandafter\def\csname BookIndex@T0433\endcsname{jieheduomianti@结合多面体（associahedron）}
\expandafter\let\csname BookOriginal@trAssociahedron\endcsname\trAssociahedron
\renewcommand{\trAssociahedron}{\BookTerm{T0433}{\csname BookOriginal@trAssociahedron\endcsname}}
\expandafter\def\csname BookIndex@T0434\endcsname{wendingquxian@稳定曲线（stable curve）}
\expandafter\let\csname BookOriginal@trStableCurve\endcsname\trStableCurve
\renewcommand{\trStableCurve}{\BookTerm{T0434}{\csname BookOriginal@trStableCurve\endcsname}}
\expandafter\def\csname BookIndex@T0435\endcsname{shiguiji@实轨迹（real locus）}
\expandafter\let\csname BookOriginal@trRealLocus\endcsname\trRealLocus
\renewcommand{\trRealLocus}{\BookTerm{T0435}{\csname BookOriginal@trRealLocus\endcsname}}
\expandafter\def\csname BookIndex@T0436\endcsname{jiedian@结点（node）}
\expandafter\let\csname BookOriginal@trNode\endcsname\trNode
\renewcommand{\trNode}{\BookTerm{T0436}{\csname BookOriginal@trNode\endcsname}}
\expandafter\def\csname BookIndex@T0437\endcsname{fanquanchun@反全纯（antiholomorphic）}
\expandafter\let\csname BookOriginal@trAntiholomorphic\endcsname\trAntiholomorphic
\renewcommand{\trAntiholomorphic}{\BookTerm{T0437}{\csname BookOriginal@trAntiholomorphic\endcsname}}
\expandafter\def\csname BookIndex@T0438\endcsname{liantonghe@连通和（connected sum）}
\expandafter\let\csname BookOriginal@trConnectedSum\endcsname\trConnectedSum
\renewcommand{\trConnectedSum}{\BookTerm{T0438}{\csname BookOriginal@trConnectedSum\endcsname}}
\expandafter\def\csname BookIndex@T0439\endcsname{waibuqi@外部旗（exterior flag）}
\expandafter\let\csname BookOriginal@trExteriorFlag\endcsname\trExteriorFlag
\renewcommand{\trExteriorFlag}{\BookTerm{T0439}{\csname BookOriginal@trExteriorFlag\endcsname}}
\expandafter\def\csname BookIndex@T0440\endcsname{houbao@厚--薄（thick thin）}
\expandafter\let\csname BookOriginal@trThickThin\endcsname\trThickThin
\renewcommand{\trThickThin}{\BookTerm{T0440}{\csname BookOriginal@trThickThin\endcsname}}
\expandafter\def\csname BookIndex@T0441\endcsname{tasuo@塌缩（collapse）}
\expandafter\let\csname BookOriginal@trCollapse\endcsname\trCollapse
\renewcommand{\trCollapse}{\BookTerm{T0441}{\csname BookOriginal@trCollapse\endcsname}}
\expandafter\def\csname BookIndex@T0442\endcsname{chuankong@穿孔（punctured）}
\expandafter\let\csname BookOriginal@trPunctured\endcsname\trPunctured
\renewcommand{\trPunctured}{\BookTerm{T0442}{\csname BookOriginal@trPunctured\endcsname}}
\expandafter\def\csname BookIndex@T0443\endcsname{duioucong@对偶丛（dual bundle）}
\expandafter\let\csname BookOriginal@trDualBundle\endcsname\trDualBundle
\renewcommand{\trDualBundle}{\BookTerm{T0443}{\csname BookOriginal@trDualBundle\endcsname}}
\expandafter\def\csname BookIndex@T0444\endcsname{bianjiefenzhi@边界分支（boundary component）}
\expandafter\let\csname BookOriginal@trBoundaryComponent\endcsname\trBoundaryComponent
\renewcommand{\trBoundaryComponent}{\BookTerm{T0444}{\csname BookOriginal@trBoundaryComponent\endcsname}}
\expandafter\def\csname BookIndex@T0445\endcsname{jidashu@极大树（maximal tree）}
\expandafter\let\csname BookOriginal@trMaximalTree\endcsname\trMaximalTree
\renewcommand{\trMaximalTree}{\BookTerm{T0445}{\csname BookOriginal@trMaximalTree\endcsname}}
\expandafter\def\csname BookIndex@T0446\endcsname{duozhongfugai@多重覆盖（multiple cover）}
\expandafter\let\csname BookOriginal@trMultipleCover\endcsname\trMultipleCover
\renewcommand{\trMultipleCover}{\BookTerm{T0446}{\csname BookOriginal@trMultipleCover\endcsname}}
\expandafter\def\csname BookIndex@T0447\endcsname{wufenci@无分次（ungraded）}
\expandafter\let\csname BookOriginal@trUngraded\endcsname\trUngraded
\renewcommand{\trUngraded}{\BookTerm{T0447}{\csname BookOriginal@trUngraded\endcsname}}
\expandafter\def\csname BookIndex@T0448\endcsname{weiyiyanta@唯一延拓（unique continuation）}
\expandafter\let\csname BookOriginal@trUniqueContinuation\endcsname\trUniqueContinuation
\renewcommand{\trUniqueContinuation}{\BookTerm{T0448}{\csname BookOriginal@trUniqueContinuation\endcsname}}
\expandafter\def\csname BookIndex@T0449\endcsname{liangdingyixing@良定义性（well definedness）}
\expandafter\let\csname BookOriginal@trWellDefinedness\endcsname\trWellDefinedness
\renewcommand{\trWellDefinedness}{\BookTerm{T0449}{\csname BookOriginal@trWellDefinedness\endcsname}}
\expandafter\def\csname BookIndex@T0450\endcsname{rongguan@融贯（coherent）}
\expandafter\let\csname BookOriginal@trCoherent\endcsname\trCoherent
\renewcommand{\trCoherent}{\BookTerm{T0450}{\csname BookOriginal@trCoherent\endcsname}}
\expandafter\def\csname BookIndex@T0451\endcsname{quanzifanchou@全子范畴（full subcategory）}
\expandafter\let\csname BookOriginal@trFullSubcategory\endcsname\trFullSubcategory
\renewcommand{\trFullSubcategory}{\BookTerm{T0451}{\csname BookOriginal@trFullSubcategory\endcsname}}
\expandafter\def\csname BookIndex@T0452\endcsname{zifanchou@子范畴（subcategory）}
\expandafter\let\csname BookOriginal@trSubcategory\endcsname\trSubcategory
\renewcommand{\trSubcategory}{\BookTerm{T0452}{\csname BookOriginal@trSubcategory\endcsname}}
\expandafter\def\csname BookIndex@T0453\endcsname{yanchihanshu@延迟函数（delay function）}
\expandafter\let\csname BookOriginal@trDelayFunction\endcsname\trDelayFunction
\renewcommand{\trDelayFunction}{\BookTerm{T0453}{\csname BookOriginal@trDelayFunction\endcsname}}
\expandafter\def\csname BookIndex@T0454\endcsname{zidengjia@自等价（autoequivalence）}
\expandafter\let\csname BookOriginal@trAutoequivalence\endcsname\trAutoequivalence
\renewcommand{\trAutoequivalence}{\BookTerm{T0454}{\csname BookOriginal@trAutoequivalence\endcsname}}
\expandafter\def\csname BookIndex@T0455\endcsname{qunpei@群胚（groupoid）}
\expandafter\let\csname BookOriginal@trGroupoid\endcsname\trGroupoid
\renewcommand{\trGroupoid}{\BookTerm{T0455}{\csname BookOriginal@trGroupoid\endcsname}}
\expandafter\def\csname BookIndex@T0456\endcsname{xiancong@线丛（line bundle）}
\expandafter\let\csname BookOriginal@trLineBundle\endcsname\trLineBundle
\renewcommand{\trLineBundle}{\BookTerm{T0456}{\csname BookOriginal@trLineBundle\endcsname}}
\expandafter\def\csname BookIndex@T0457\endcsname{waimi@外幂（exterior power）}
\expandafter\let\csname BookOriginal@trExteriorPower\endcsname\trExteriorPower
\renewcommand{\trExteriorPower}{\BookTerm{T0457}{\csname BookOriginal@trExteriorPower\endcsname}}
\expandafter\def\csname BookIndex@T0458\endcsname{xiangjiaoshu@相交数（intersection number）}
\expandafter\let\csname BookOriginal@trIntersectionNumber\endcsname\trIntersectionNumber
\renewcommand{\trIntersectionNumber}{\BookTerm{T0458}{\csname BookOriginal@trIntersectionNumber\endcsname}}
\expandafter\def\csname BookIndex@T0459\endcsname{fencha@分岔（bifurcation）}
\expandafter\let\csname BookOriginal@trBifurcation\endcsname\trBifurcation
\renewcommand{\trBifurcation}{\BookTerm{T0459}{\csname BookOriginal@trBifurcation\endcsname}}
\expandafter\def\csname BookIndex@T0460\endcsname{yuzhi@余秩（corank）}
\expandafter\let\csname BookOriginal@trCorank\endcsname\trCorank
\renewcommand{\trCorank}{\BookTerm{T0460}{\csname BookOriginal@trCorank\endcsname}}
\expandafter\def\csname BookIndex@T0461\endcsname{lingkongjian@零空间（null space）}
\expandafter\let\csname BookOriginal@trNullSpace\endcsname\trNullSpace
\renewcommand{\trNullSpace}{\BookTerm{T0461}{\csname BookOriginal@trNullSpace\endcsname}}
\expandafter\def\csname BookIndex@T0462\endcsname{lianhuan@链环（link）}
\expandafter\let\csname BookOriginal@trLink\endcsname\trLink
\renewcommand{\trLink}{\BookTerm{T0462}{\csname BookOriginal@trLink\endcsname}}
\expandafter\def\csname BookIndex@T0463\endcsname{wendinghua@稳定化（stabilization）}
\expandafter\let\csname BookOriginal@trStabilization\endcsname\trStabilization
\renewcommand{\trStabilization}{\BookTerm{T0463}{\csname BookOriginal@trStabilization\endcsname}}
\expandafter\def\csname BookIndex@T0464\endcsname{tezhengxiangliang@特征向量（eigenvector）}
\expandafter\let\csname BookOriginal@trEigenvector\endcsname\trEigenvector
\renewcommand{\trEigenvector}{\BookTerm{T0464}{\csname BookOriginal@trEigenvector\endcsname}}
\expandafter\def\csname BookIndex@T0465\endcsname{jiaochaxingshi@交叉形式（crossing form）}
\expandafter\let\csname BookOriginal@trCrossingForm\endcsname\trCrossingForm
\renewcommand{\trCrossingForm}{\BookTerm{T0465}{\csname BookOriginal@trCrossingForm\endcsname}}
\expandafter\def\csname BookIndex@T0466\endcsname{mojiegou@膜结构（brane structure）}
\expandafter\let\csname BookOriginal@trBraneStructure\endcsname\trBraneStructure
\renewcommand{\trBraneStructure}{\BookTerm{T0466}{\csname BookOriginal@trBraneStructure\endcsname}}
\expandafter\def\csname BookIndex@T0467\endcsname{zhongxinkuozhang@中心扩张（central extension）}
\expandafter\let\csname BookOriginal@trCentralExtension\endcsname\trCentralExtension
\renewcommand{\trCentralExtension}{\BookTerm{T0467}{\csname BookOriginal@trCentralExtension\endcsname}}
\expandafter\def\csname BookIndex@T0468\endcsname{zhucong@主丛（principal bundle）}
\expandafter\let\csname BookOriginal@trPrincipalBundle\endcsname\trPrincipalBundle
\renewcommand{\trPrincipalBundle}{\BookTerm{T0468}{\csname BookOriginal@trPrincipalBundle\endcsname}}
\expandafter\def\csname BookIndex@T0469\endcsname{shixiancong@实线丛（real line bundle）}
\expandafter\let\csname BookOriginal@trRealLineBundle\endcsname\trRealLineBundle
\renewcommand{\trRealLineBundle}{\BookTerm{T0469}{\csname BookOriginal@trRealLineBundle\endcsname}}
\expandafter\def\csname BookIndex@T0470\endcsname{erzhongfugai@二重覆盖（double cover）}
\expandafter\let\csname BookOriginal@trDoubleCover\endcsname\trDoubleCover
\renewcommand{\trDoubleCover}{\BookTerm{T0470}{\csname BookOriginal@trDoubleCover\endcsname}}
\expandafter\def\csname BookIndex@T0471\endcsname{zhuqixingkongjian@主齐性空间（principal homogeneous space）}
\expandafter\let\csname BookOriginal@trPrincipalHomogeneousSpace\endcsname\trPrincipalHomogeneousSpace
\renewcommand{\trPrincipalHomogeneousSpace}{\BookTerm{T0471}{\csname BookOriginal@trPrincipalHomogeneousSpace\endcsname}}
\expandafter\def\csname BookIndex@T0472\endcsname{fanshe@反射（reflection）}
\expandafter\let\csname BookOriginal@trReflection\endcsname\trReflection
\renewcommand{\trReflection}{\BookTerm{T0472}{\csname BookOriginal@trReflection\endcsname}}
\expandafter\def\csname BookIndex@T0473\endcsname{dingxiangfanzhuan@定向反转（orientation reversing）}
\expandafter\let\csname BookOriginal@trOrientationReversing\endcsname\trOrientationReversing
\renewcommand{\trOrientationReversing}{\BookTerm{T0473}{\csname BookOriginal@trOrientationReversing\endcsname}}
\expandafter\def\csname BookIndex@T0474\endcsname{niesuo@捏缩（pinch）}
\expandafter\let\csname BookOriginal@trPinch\endcsname\trPinch
\renewcommand{\trPinch}{\BookTerm{T0474}{\csname BookOriginal@trPinch\endcsname}}
\expandafter\def\csname BookIndex@T0475\endcsname{anuodefenceng@阿诺德分层（Arnold stratification）}
\expandafter\let\csname BookOriginal@trArnoldStratification\endcsname\trArnoldStratification
\renewcommand{\trArnoldStratification}{\BookTerm{T0475}{\csname BookOriginal@trArnoldStratification\endcsname}}
\expandafter\def\csname BookIndex@T0476\endcsname{pingfangxiangwei@平方相位（squared phase）}
\expandafter\let\csname BookOriginal@trSquaredPhase\endcsname\trSquaredPhase
\renewcommand{\trSquaredPhase}{\BookTerm{T0476}{\csname BookOriginal@trSquaredPhase\endcsname}}
\expandafter\def\csname BookIndex@T0477\endcsname{shizhixiangwei@实值相位（real valued phase）}
\expandafter\let\csname BookOriginal@trRealValuedPhase\endcsname\trRealValuedPhase
\renewcommand{\trRealValuedPhase}{\BookTerm{T0477}{\csname BookOriginal@trRealValuedPhase\endcsname}}
\expandafter\def\csname BookIndex@T0478\endcsname{dingxiangsuanzi@定向算子（orientation operator）}
\expandafter\let\csname BookOriginal@trOrientationOperator\endcsname\trOrientationOperator
\renewcommand{\trOrientationOperator}{\BookTerm{T0478}{\csname BookOriginal@trOrientationOperator\endcsname}}
\expandafter\def\csname BookIndex@T0479\endcsname{guifanhua@规范化（normalization）}
\expandafter\let\csname BookOriginal@trNormalization\endcsname\trNormalization
\renewcommand{\trNormalization}{\BookTerm{T0479}{\csname BookOriginal@trNormalization\endcsname}}
\expandafter\def\csname BookIndex@T0480\endcsname{gelasimanliuxing@格拉斯曼流形（grassmannian）}
\expandafter\let\csname BookOriginal@trGrassmannian\endcsname\trGrassmannian
\renewcommand{\trGrassmannian}{\BookTerm{T0480}{\csname BookOriginal@trGrassmannian\endcsname}}
\expandafter\def\csname BookIndex@T0481\endcsname{bianjieyuanzhou@边界圆周（boundary circle）}
\expandafter\let\csname BookOriginal@trBoundaryCircle\endcsname\trBoundaryCircle
\renewcommand{\trBoundaryCircle}{\BookTerm{T0481}{\csname BookOriginal@trBoundaryCircle\endcsname}}
\expandafter\def\csname BookIndex@T0482\endcsname{yingzi@影子（shadow）}
\expandafter\let\csname BookOriginal@trShadow\endcsname\trShadow
\renewcommand{\trShadow}{\BookTerm{T0482}{\csname BookOriginal@trShadow\endcsname}}
\expandafter\def\csname BookIndex@T0483\endcsname{faxiangjiao@法向角（normal angle）}
\expandafter\let\csname BookOriginal@trNormalAngle\endcsname\trNormalAngle
\renewcommand{\trNormalAngle}{\BookTerm{T0483}{\csname BookOriginal@trNormalAngle\endcsname}}
\expandafter\def\csname BookIndex@T0484\endcsname{bufenkejixing@部分可积性（partial integrability）}
\expandafter\let\csname BookOriginal@trPartialIntegrability\endcsname\trPartialIntegrability
\renewcommand{\trPartialIntegrability}{\BookTerm{T0484}{\csname BookOriginal@trPartialIntegrability\endcsname}}
\expandafter\def\csname BookIndex@T0485\endcsname{zidongzhengzexing@自动正则性（automatic regularity）}
\expandafter\let\csname BookOriginal@trAutomaticRegularity\endcsname\trAutomaticRegularity
\renewcommand{\trAutomaticRegularity}{\BookTerm{T0485}{\csname BookOriginal@trAutomaticRegularity\endcsname}}
\expandafter\def\csname BookIndex@T0486\endcsname{zuhefuluoershangtongdiao@组合弗洛尔上同调（combinatorial Floer cohomology）}
\expandafter\let\csname BookOriginal@trCombinatorialFloerCohomology\endcsname\trCombinatorialFloerCohomology
\renewcommand{\trCombinatorialFloerCohomology}{\BookTerm{T0486}{\csname BookOriginal@trCombinatorialFloerCohomology\endcsname}}
\expandafter\def\csname BookIndex@T0487\endcsname{yezhuangjiegou@叶状结构（foliation）}
\expandafter\let\csname BookOriginal@trFoliation\endcsname\trFoliation
\renewcommand{\trFoliation}{\BookTerm{T0487}{\csname BookOriginal@trFoliation\endcsname}}
\expandafter\def\csname BookIndex@T0488\endcsname{guaiqiumian@怪球面（exotic sphere）}
\expandafter\let\csname BookOriginal@trExoticSphere\endcsname\trExoticSphere
\renewcommand{\trExoticSphere}{\BookTerm{T0488}{\csname BookOriginal@trExoticSphere\endcsname}}
\expandafter\def\csname BookIndex@T0489\endcsname{shuangxianweihua@双纤维化（bifibration）}
\expandafter\let\csname BookOriginal@trBifibration\endcsname\trBifibration
\renewcommand{\trBifibration}{\BookTerm{T0489}{\csname BookOriginal@trBifibration\endcsname}}
\expandafter\def\csname BookIndex@T0490\endcsname{jialinjiezhi@假临界值（fake critical value）}
\expandafter\let\csname BookOriginal@trFakeCriticalValue\endcsname\trFakeCriticalValue
\renewcommand{\trFakeCriticalValue}{\BookTerm{T0490}{\csname BookOriginal@trFakeCriticalValue\endcsname}}
\expandafter\def\csname BookIndex@T0491\endcsname{daiqidianxianweihua@带奇点纤维化（fibration with singularities）}
\expandafter\let\csname BookOriginal@trFibrationWithSingularities\endcsname\trFibrationWithSingularities
\renewcommand{\trFibrationWithSingularities}{\BookTerm{T0491}{\csname BookOriginal@trFibrationWithSingularities\endcsname}}
\expandafter\def\csname BookIndex@T0492\endcsname{liangtai@良态（well behaved）}
\expandafter\let\csname BookOriginal@trWellBehaved\endcsname\trWellBehaved
\renewcommand{\trWellBehaved}{\BookTerm{T0492}{\csname BookOriginal@trWellBehaved\endcsname}}
\expandafter\def\csname BookIndex@T0493\endcsname{zhuxianweichangzhi@逐纤维常值（constant along fibres）}
\expandafter\let\csname BookOriginal@trConstantAlongFibres\endcsname\trConstantAlongFibres
\renewcommand{\trConstantAlongFibres}{\BookTerm{T0493}{\csname BookOriginal@trConstantAlongFibres\endcsname}}
\expandafter\def\csname BookIndex@T0494\endcsname{banniuzhuan@半扭转（half twist）}
\expandafter\let\csname BookOriginal@trHalfTwist\endcsname\trHalfTwist
\renewcommand{\trHalfTwist}{\BookTerm{T0494}{\csname BookOriginal@trHalfTwist\endcsname}}
\expandafter\def\csname BookIndex@T0495\endcsname{teyiji@特异基（distinguished basis）}
\expandafter\let\csname BookOriginal@trDistinguishedBasis\endcsname\trDistinguishedBasis
\renewcommand{\trDistinguishedBasis}{\BookTerm{T0495}{\csname BookOriginal@trDistinguishedBasis\endcsname}}
\expandafter\def\csname BookIndex@T0496\endcsname{jilujing@基路径（base path）}
\expandafter\let\csname BookOriginal@trBasePath\endcsname\trBasePath
\renewcommand{\trBasePath}{\BookTerm{T0496}{\csname BookOriginal@trBasePath\endcsname}}
\expandafter\def\csname BookIndex@T0497\endcsname{pipeitonghen@匹配同痕（matching isotopy）}
\expandafter\let\csname BookOriginal@trMatchingIsotopy\endcsname\trMatchingIsotopy
\renewcommand{\trMatchingIsotopy}{\BookTerm{T0497}{\csname BookOriginal@trMatchingIsotopy\endcsname}}
\expandafter\def\csname BookIndex@T0498\endcsname{jishujiemian@计数截面（counting sections）}
\expandafter\let\csname BookOriginal@trCountingSections\endcsname\trCountingSections
\renewcommand{\trCountingSections}{\BookTerm{T0498}{\csname BookOriginal@trCountingSections\endcsname}}
\expandafter\def\csname BookIndex@T0499\endcsname{daidaizhuangduandexianweihua@带带状端的纤维化（fibration with striplike ends）}
\expandafter\let\csname BookOriginal@trFibrationWithStriplikeEnds\endcsname\trFibrationWithStriplikeEnds
\renewcommand{\trFibrationWithStriplikeEnds}{\BookTerm{T0499}{\csname BookOriginal@trFibrationWithStriplikeEnds\endcsname}}
\expandafter\def\csname BookIndex@T0500\endcsname{xiaomeidingli@消没定理（vanishing theorem）}
\expandafter\let\csname BookOriginal@trVanishingTheorem\endcsname\trVanishingTheorem
\renewcommand{\trVanishingTheorem}{\BookTerm{T0500}{\csname BookOriginal@trVanishingTheorem\endcsname}}
\expandafter\def\csname BookIndex@T0501\endcsname{jibentaishe@基本态射（basic morphism）}
\expandafter\let\csname BookOriginal@trBasicMorphism\endcsname\trBasicMorphism
\renewcommand{\trBasicMorphism}{\BookTerm{T0501}{\csname BookOriginal@trBasicMorphism\endcsname}}
\expandafter\def\csname BookIndex@T0502\endcsname{fuzhulianfuxing@辅助链复形（auxiliary chain complex）}
\expandafter\let\csname BookOriginal@trAuxiliaryChainComplex\endcsname\trAuxiliaryChainComplex
\renewcommand{\trAuxiliaryChainComplex}{\BookTerm{T0502}{\csname BookOriginal@trAuxiliaryChainComplex\endcsname}}
\expandafter\def\csname BookIndex@T0503\endcsname{daishuniuzhuan@代数扭转（algebraic twisting）}
\expandafter\let\csname BookOriginal@trAlgebraicTwisting\endcsname\trAlgebraicTwisting
\renewcommand{\trAlgebraicTwisting}{\BookTerm{T0503}{\csname BookOriginal@trAlgebraicTwisting\endcsname}}
\expandafter\def\csname BookIndex@T0504\endcsname{jiheniuzhuan@几何扭转（geometric twisting）}
\expandafter\let\csname BookOriginal@trGeometricTwisting\endcsname\trGeometricTwisting
\renewcommand{\trGeometricTwisting}{\BookTerm{T0504}{\csname BookOriginal@trGeometricTwisting\endcsname}}
\expandafter\def\csname BookIndex@T0505\endcsname{nidengjia@拟等价（quasiequivalent）}
\expandafter\let\csname BookOriginal@trQuasiequivalent\endcsname\trQuasiequivalent
\renewcommand{\trQuasiequivalent}{\BookTerm{T0505}{\csname BookOriginal@trQuasiequivalent\endcsname}}
\expandafter\def\csname BookIndex@T0506\endcsname{erzhongfenqifugai@二重分歧覆盖（double branched cover）}
\expandafter\let\csname BookOriginal@trDoubleBranchedCover\endcsname\trDoubleBranchedCover
\renewcommand{\trDoubleBranchedCover}{\BookTerm{T0506}{\csname BookOriginal@trDoubleBranchedCover\endcsname}}
\expandafter\def\csname BookIndex@T0507\endcsname{jiantu@箭图（quiver）}
\expandafter\let\csname BookOriginal@trQuiver\endcsname\trQuiver
\renewcommand{\trQuiver}{\BookTerm{T0507}{\csname BookOriginal@trQuiver\endcsname}}
\expandafter\def\csname BookIndex@T0508\endcsname{lujingdaishu@路径代数（path algebra）}
\expandafter\let\csname BookOriginal@trPathAlgebra\endcsname\trPathAlgebra
\renewcommand{\trPathAlgebra}{\BookTerm{T0508}{\csname BookOriginal@trPathAlgebra\endcsname}}
\expandafter\def\csname BookIndex@T0509\endcsname{duocidiaohe@多次调和（plurisubharmonic）}
\expandafter\let\csname BookOriginal@trPlurisubharmonic\endcsname\trPlurisubharmonic
\renewcommand{\trPlurisubharmonic}{\BookTerm{T0509}{\csname BookOriginal@trPlurisubharmonic\endcsname}}
\expandafter\def\csname BookIndex@T0510\endcsname{qiongjie@穷竭（exhausting）}
\expandafter\let\csname BookOriginal@trExhausting\endcsname\trExhausting
\renewcommand{\trExhausting}{\BookTerm{T0510}{\csname BookOriginal@trExhausting\endcsname}}
\expandafter\def\csname BookIndex@T0511\endcsname{youxianxing@有限型（finite type）}
\expandafter\let\csname BookOriginal@trFiniteType\endcsname\trFiniteType
\renewcommand{\trFiniteType}{\BookTerm{T0511}{\csname BookOriginal@trFiniteType\endcsname}}
\expandafter\def\csname BookIndex@T0512\endcsname{chuzi@除子（divisor）}
\expandafter\let\csname BookOriginal@trDivisor\endcsname\trDivisor
\renewcommand{\trDivisor}{\BookTerm{T0512}{\csname BookOriginal@trDivisor\endcsname}}
\expandafter\def\csname BookIndex@T0513\endcsname{zhengguijiaocha@正规交叉（normal crossings）}
\expandafter\let\csname BookOriginal@trNormalCrossings\endcsname\trNormalCrossings
\renewcommand{\trNormalCrossings}{\BookTerm{T0513}{\csname BookOriginal@trNormalCrossings\endcsname}}
\expandafter\def\csname BookIndex@T0514\endcsname{jiguiji@基轨迹（base locus）}
\expandafter\let\csname BookOriginal@trBaseLocus\endcsname\trBaseLocus
\renewcommand{\trBaseLocus}{\BookTerm{T0514}{\csname BookOriginal@trBaseLocus\endcsname}}
\expandafter\def\csname BookIndex@T0515\endcsname{budingguiji@不定轨迹（indeterminacy locus）}
\expandafter\let\csname BookOriginal@trIndeterminacyLocus\endcsname\trIndeterminacyLocus
\renewcommand{\trIndeterminacyLocus}{\BookTerm{T0515}{\csname BookOriginal@trIndeterminacyLocus\endcsname}}
\expandafter\def\csname BookIndex@T0516\endcsname{liwaichuzi@例外除子（exceptional divisor）}
\expandafter\let\csname BookOriginal@trExceptionalDivisor\endcsname\trExceptionalDivisor
\renewcommand{\trExceptionalDivisor}{\BookTerm{T0516}{\csname BookOriginal@trExceptionalDivisor\endcsname}}
\expandafter\def\csname BookIndex@T0517\endcsname{baopo@爆破（blowup）}
\expandafter\let\csname BookOriginal@trBlowup\endcsname\trBlowup
\renewcommand{\trBlowup}{\BookTerm{T0517}{\csname BookOriginal@trBlowup\endcsname}}
\expandafter\def\csname BookIndex@T0518\endcsname{duoqiuti@多球体（polyball）}
\expandafter\let\csname BookOriginal@trPolyball\endcsname\trPolyball
\renewcommand{\trPolyball}{\BookTerm{T0518}{\csname BookOriginal@trPolyball\endcsname}}
\expandafter\def\csname BookIndex@T0519\endcsname{kaileishi@凯勒势（Kähler potential）}
\expandafter\let\csname BookOriginal@trKaehlerPotential\endcsname\trKaehlerPotential
\renewcommand{\trKaehlerPotential}{\BookTerm{T0519}{\csname BookOriginal@trKaehlerPotential\endcsname}}
\expandafter\def\csname BookIndex@T0520\endcsname{chaopingmian@超平面（hyperplane）}
\expandafter\let\csname BookOriginal@trHyperplane\endcsname\trHyperplane
\renewcommand{\trHyperplane}{\BookTerm{T0520}{\csname BookOriginal@trHyperplane\endcsname}}
\expandafter\def\csname BookIndex@T0521\endcsname{yibanweizhi@一般位置（general position）}
\expandafter\let\csname BookOriginal@trGeneralPosition\endcsname\trGeneralPosition
\renewcommand{\trGeneralPosition}{\BookTerm{T0521}{\csname BookOriginal@trGeneralPosition\endcsname}}
\expandafter\def\csname BookIndex@T0522\endcsname{fazhui@法锥（normal cone）}
\expandafter\let\csname BookOriginal@trNormalCone\endcsname\trNormalCone
\renewcommand{\trNormalCone}{\BookTerm{T0522}{\csname BookOriginal@trNormalCone\endcsname}}
\expandafter\def\csname BookIndex@T0523\endcsname{xinqiege@辛切割（symplectic cut）}
\expandafter\let\csname BookOriginal@trSymplecticCut\endcsname\trSymplecticCut
\renewcommand{\trSymplecticCut}{\BookTerm{T0523}{\csname BookOriginal@trSymplecticCut\endcsname}}
\expandafter\def\csname BookIndex@T0524\endcsname{xuanzhuanshu@旋转数（rotation number）}
\expandafter\let\csname BookOriginal@trRotationNumber\endcsname\trRotationNumber
\renewcommand{\trRotationNumber}{\BookTerm{T0524}{\csname BookOriginal@trRotationNumber\endcsname}}
\expandafter\def\csname BookIndex@T0525\endcsname{yingsheleiqun@映射类群（mapping class group）}
\expandafter\let\csname BookOriginal@trMappingClassGroup\endcsname\trMappingClassGroup
\renewcommand{\trMappingClassGroup}{\BookTerm{T0525}{\csname BookOriginal@trMappingClassGroup\endcsname}}
\expandafter\def\csname BookIndex@T0526\endcsname{erciduiou@二次对偶（quadratic dual）}
\expandafter\let\csname BookOriginal@trQuadraticDual\endcsname\trQuadraticDual
\renewcommand{\trQuadraticDual}{\BookTerm{T0526}{\csname BookOriginal@trQuadraticDual\endcsname}}
\makeatother
